% Section file for Chapter: Twisted Holography and Quantum Gravity
% Result (G10): Fredholm Partition Functions
%
% This section develops the nonperturbative partition function
% as a Fredholm determinant on the one-particle Bergman space,
% extending from the Heisenberg prototype through the full
% standard landscape.

% ==========================================
% LOCAL MACRO PROVISIONS
% ==========================================
\providecommand{\Bergman}{\mathcal{B}}
\providecommand{\Fock}{\mathcal{F}}
\providecommand{\Hilb}{\mathcal{H}\!\mathit{ilb}}
\providecommand{\sewop}{T}
\providecommand{\sewker}{K}
\providecommand{\trcl}{\mathcal{L}^1}
\providecommand{\hscl}{\mathcal{L}^2}
\providecommand{\freddet}{\mathrm{det}_{\mathrm{Fred}}}
\providecommand{\Wick}{\mathcal{W}}
\providecommand{\Bergker}{\mathsf{B}}
\providecommand{\primeform}{E}
\providecommand{\collar}{\mathfrak{c}}
\providecommand{\sewenv}{\cA^{\mathrm{sew}}}
\providecommand{\anbar}{\barB^{\mathrm{an}}}
\providecommand{\algcore}{\cA_{\mathrm{alg}}}
\providecommand{\codercat}{D^{\mathrm{co}}}
\providecommand{\contrcat}{D^{\mathrm{ctr}}}
\providecommand{\graphnorm}{\|\cdot\|_{\mathrm{graph}}}
\providecommand{\secquant}{\Gamma}
\providecommand{\hodgebdle}{\mathbb{E}}
\providecommand{\hodgedet}{\det\hodgebdle}
\providecommand{\zetareg}{\det\nolimits'_\zeta}
\providecommand{\dedeta}{\eta}
\providecommand{\Sp}{\mathrm{Sp}}
\providecommand{\Siegel}{\mathfrak{H}}
\providecommand{\OPEcoeff}{C}

\section{Fredholm partition functions}
% label removed: sec:thqg-fredholm-partition-functions
\index{Fredholm determinant!partition function|textbf}
\index{partition function!Fredholm}

The preceding nine results of this chapter extract holographic,
gravitational, and bootstrap consequences from the universal
Maurer--Cartan element
$\Theta_\cA \in \MC(\gAmod)$
at the level of formal algebraic identities.  This final result
crosses from the algebraic to the analytic: the partition functions
are convergent Fredholm determinants on explicit Hilbert spaces.

The key bridge is the \emph{sewing envelope}
$\sewenv$
(Definition~\ref{def:sewing-envelope}): the Hausdorff completion
of the algebraic chiral core $\algcore$ for the locally convex
topology generated by all amplitude seminorms.  For the Heisenberg
algebra $\cH_\kappa$, the sewing envelope is the symmetric algebra
of the Bergman space $A^2(D)$ of the disk
(Theorem~\ref{thm:heisenberg-sewing}\,(i)), and the genus-$g$
partition function is the Fredholm determinant of an explicit
trace-class sewing operator on $A^2(D)$.  The HS-sewing criterion
(Theorem~\ref{thm:general-hs-sewing}) extends the convergence to
the entire standard landscape of universal algebras (convergence,
not Fredholm
form).  On the proved Gaussian scalar lane, the scalar package is
a $\kappa$-multiple of the Heisenberg one, so the Fredholm
determinant structure propagates there.  For classes L, C, and M,
the partition function has additional non-Fredholm corrections
from the higher-arity shadow obstruction tower
(\S\ref{subsec:thqg-X-non-fredholm}).

The section is organized as follows.
\begin{enumerate}[label=\textup{(\arabic*)}]
\item Sewing envelopes and Bergman spaces:
  the analytic framework for one-particle reduction
  (\S\ref{subsec:thqg-X-sewing-bergman}).
\item The Heisenberg sewing theorem in full:
  four clauses with complete proofs
  (\S\ref{subsec:thqg-X-heisenberg-sewing}).
\item One-particle reduction, trace-class estimates,
  and the Dedekind eta derivation
  (\S\ref{subsec:thqg-X-one-particle}).
\item Extension to class-G algebras via scalar saturation
  (\S\ref{subsec:thqg-X-class-G}).
\item Non-Fredholm corrections for classes L, C, M
  (\S\ref{subsec:thqg-X-non-fredholm}).
\item Analytic aspects: the analytic bar coalgebra,
  coderived categories, and open problems
  (\S\ref{subsec:thqg-X-analytic-aspects}).
\end{enumerate}

% =====================================================================
%
%  SECTION 1: SEWING ENVELOPES AND BERGMAN SPACES
%
% =====================================================================

\subsection{Sewing envelopes and Bergman spaces}
% label removed: subsec:thqg-X-sewing-bergman

\subsubsection{The Bergman space of a disk}

\begin{definition}[Bergman space of a disk]%
% label removed: def:thqg-X-bergman-space%
\index{Bergman space|textbf}%
For $D = \{z \in \C : |z| < 1\}$ the \textbf{unit disk},
the \emph{Bergman space} $A^2(D)$ is the Hilbert space of
holomorphic functions that are square-integrable with respect
to the area measure:
\begin{equation}% label removed: eq:thqg-X-bergman-norm
A^2(D)
\;:=\;
\Bigl\{f \in \cO(D) : \|f\|^2_{A^2}
= \int_D |f(z)|^2\,dA(z) < \infty\Bigr\},
\end{equation}
where $dA(z) = \frac{i}{2}\,dz \wedge d\bar{z}
= r\,dr\,d\theta$ in polar coordinates.  The monomials
$e_n(z) = \sqrt{(n+1)/\pi}\, z^n$ ($n \geq 0$)
form an orthonormal basis:
$\langle e_n, e_m\rangle = \delta_{nm}$ (verified by
$\int_D |z|^{2n}\,dA = \pi/(n+1)$).  The
\emph{reproducing kernel} (Bergman kernel for the unit
disk) is
\begin{equation}% label removed: eq:thqg-X-bergman-kernel
\Bergker_D(z,w)
\;=\;
\sum_{n=0}^{\infty} e_n(z)\overline{e_n(w)}
\;=\;
\frac{1}{\pi}\sum_{n=0}^{\infty}(n+1)(z\bar{w})^n
\;=\;
\frac{1}{\pi(1 - z\bar{w})^2}.
\end{equation}
The last identity uses $\sum_{n=0}^{\infty}(n+1)x^n
= 1/(1-x)^2$ for $|x| < 1$.  \emph{This formula is
specific to the unit disk}; for a general bounded domain
$\Omega \subset \C$, the Bergman kernel depends on the
conformal geometry of~$\Omega$.
\end{definition}

\begin{remark}[Conformal weight and sections of $K^{-1/2}$]%
% label removed: rem:thqg-X-bergman-conformal%
The prime form $\primeform(z_1,z_2)$ on a Riemann surface is a section
of $K^{-1/2} \boxtimes K^{-1/2}$, \emph{not} $K^{+1/2}$.
The Bergman kernel
$\Bergker_D(z,w)$ is a section of $K \boxtimes \bar{K}$, consistent
with its conformal weight $(1,1)$.  The mode--Bergman correspondence
(Remark~\ref{rem:heisenberg-mode-bergman}) maps Fock modes to
Bergman monomials at weight $0$, absorbing the conformal weight
into the $\sqrt{n+1}$ prefactor.
\end{remark}

\begin{definition}[Bergman space with collar parameter]%
% label removed: def:thqg-X-bergman-collar%
\index{Bergman space!collar parameter}%
For $0 < q < 1$, define the \emph{weighted Bergman space}
\begin{equation}% label removed: eq:thqg-X-bergman-weighted
A^2_q(D)
\;:=\;
\Bigl\{f \in \cO(D) :
\|f\|^2_{q} := \sum_{n=0}^{\infty}
q^n \,|a_n|^2 \,(n+1)^{-1} < \infty
\;\text{where}\; f = \sum_{n=0}^{\infty} a_n z^n\Bigr\}.
\end{equation}
The collar parameter $q = e^{-t}$ ($t > 0$ the collar length)
encodes the sewing parameter for pair-of-pants decompositions.
The inclusion $A^2_{q'}(D) \hookrightarrow A^2_q(D)$ for $q' < q$
is trace class.
\end{definition}

\begin{lemma}[Bergman projection; \ClaimStatusProvedElsewhere]%
% label removed: lem:thqg-X-bergman-projection%
\index{Bergman projection}%
The orthogonal projection
$P \colon L^2(D, dA) \to A^2(D)$ is the integral operator
\[
(Pf)(z)
\;=\;
\int_D \Bergker_D(z,w)\,f(w)\,dA(w)
\;=\;
\frac{1}{\pi}\int_D \frac{f(w)}{(1-z\bar{w})^2}\,dA(w).
\]
The projection satisfies $P^2 = P$, $P^* = P$, and
$\ker P = (\bar{\partial}\text{-exact})$.
\end{lemma}

\subsubsection{The sewing envelope}

\begin{definition}[Sewing envelope; review]%
% label removed: def:thqg-X-sewing-envelope-review%
\index{sewing envelope!review}%
Recall from Definition~\ref{def:sewing-envelope} that
the \emph{sewing envelope} $\sewenv$ of an algebraic chiral
core $\algcore$ is the Hausdorff completion
for the locally convex topology generated by all amplitude
seminorms $p_{\Sigma,\xi,\eta}(a)
= |\langle \eta, A_\Sigma^{\mathrm{alg}}(\xi_1 \otimes \cdots
\otimes a \otimes \cdots \otimes \xi_m)\rangle|$,
where the supremum ranges over all finite conformally flat
bordisms $\Sigma$ with parametrized collars and all
test states $\xi, \eta$.  By Proposition~\ref{prop:sewing-universal-property},
$\sewenv$ has the universal property: every complete locally convex
realization factors uniquely through~$\sewenv$.
\end{definition}

\begin{lemma}[Amplitude seminorms form a separating family;
\ClaimStatusProvedHere]%
% label removed: lem:thqg-X-amplitude-separating%
\index{sewing envelope!separating family}%
The family of amplitude seminorms
$\{p_{\Sigma,\xi,\eta}\}$ is separating: if
$p_{\Sigma,\xi,\eta}(a) = 0$ for all bordisms $\Sigma$ with
parametrized collars and all test states $\xi, \eta$, then
$a = 0$ in~$\algcore$.
\end{lemma}

\begin{proof}
Suppose $a \in \algcore$ satisfies
$p_{\Sigma,\xi,\eta}(a) = 0$ for all $(\Sigma,\xi,\eta)$.
Then $a$ pairs trivially with all sewing amplitudes:
\[
\langle \eta,\, A_\Sigma^{\mathrm{alg}}(\xi_1 \otimes
\cdots \otimes a \otimes \cdots \otimes \xi_m)\rangle = 0
\]
for every conformally flat bordism $\Sigma$ and all test
states.  In particular, choosing $\Sigma$ to be the
pair-of-pants with collar parameter $q \to 0$
(the degeneration to the disk), the amplitude
$A_\Sigma^{\mathrm{alg}}$ reduces to the OPE
multiplication $Y(a,z)$: the condition becomes
$\langle \eta, Y(a,z)\,\xi\rangle = 0$ for all
$\xi, \eta \in \algcore$ and all $z$.  By nondegeneracy of
the OPE pairing
(Proposition~\textup{\ref{prop:sewing-universal-property}}:
the bilinear form $\langle\cdot,Y(\cdot,z)\cdot\rangle$
is nondegenerate on the algebraic core),
this forces $a = 0$.
\end{proof}

\begin{proposition}[Sewing envelope: conformal weight filtration;
\ClaimStatusProvedHere]%
% label removed: prop:thqg-X-sewing-filtration%
\index{sewing envelope!filtration}%
Let $\cA$ be a positive-energy chiral algebra with
conformal weight decomposition
$\algcore = \bigoplus_{n \geq 0} \algcore_n$, where
$\dim \algcore_n < \infty$ for all~$n$.  Then:
\begin{enumerate}[label=\textup{(\roman*)}]
\item The sewing envelope inherits a filtration
  $F^{\leq N}\sewenv
  = \overline{\bigoplus_{n \leq N} \algcore_n}$
  whose associated graded recovers the algebraic core:
  $\mathrm{gr}\,\sewenv \cong \algcore$.
\item For each $0 < q < 1$, the $q$-weighted completion
  $\sewenv_q := \widehat{\bigoplus}_{n \geq 0} q^{n/2}\,\algcore_n$
  is a Banach algebra for the operator norm of
  $A^2_q$-bounded amplitudes.
\item The sewing envelope is the projective limit
  $\sewenv = \varprojlim_{q \to 1^-} \sewenv_q$.
\end{enumerate}
\end{proposition}

\begin{proof}
Part~(i): The amplitude seminorms $p_{\Sigma,\xi,\eta}$ are
continuous with respect to the conformal weight grading because
each amplitude map preserves total conformal weight up to a shift
determined by the Euler characteristic of~$\Sigma$.
The conformal weight operator $L_0$ is bounded on each
finite-dimensional weight space $\algcore_n$ (where it acts as
scalar multiplication by~$n$), and the sewing seminorms are
dominated by $q^{L_0}$: for any bordism $\Sigma$ with collar
parameter~$q$,
\[
p_{\Sigma,\xi,\eta}(a)
\;\leq\;
C_{\Sigma,\xi,\eta}\, q^n
\quad\text{for } a \in \algcore_n.
\]
Since $q^{L_0}$ is continuous in the weight filtration
topology (it sends $F^{\leq N}\algcore$ to itself with
norm $\leq q^N$), the weight map $a \mapsto n$ is continuous
under the sewing topology.
Passage to the associated graded undoes the completion, recovering
the algebraic core.

Part~(ii): The $q$-weighting assigns norm $q^{n/2}$ to elements
of weight~$n$.  The pair-of-pants amplitude
$m \colon \algcore_a \otimes \algcore_b \to \bigoplus_c \algcore_c$
satisfies $|m^c_{a,b}| \leq K(a+b+c+1)^N$ by polynomial OPE growth.
The $q$-weighted operator norm is
$\|m\|_q \leq K \sum_c q^{c/2}(a+b+c+1)^N < \infty$
for fixed $a,b$ and $q < 1$, so $\sewenv_q$ is submultiplicative.
Completeness follows from the completeness of the
$q$-weighted $\ell^2$ space.

Part~(iii): The sewing topology is generated by
$\{p_{\Sigma,\xi,\eta}\}$, which are bounded by $q$-weighted
norms for any $q$ sufficiently close to~$1$ (depending on the
collar length of~$\Sigma$).  Hence $\sewenv$ is the intersection
$\bigcap_{q < 1} \sewenv_q = \varprojlim_{q} \sewenv_q$.
\end{proof}

\subsubsection{The Heisenberg sewing envelope}

\begin{proposition}[Heisenberg sewing envelope = $\operatorname{Sym} A^2(D)$;
\ClaimStatusProvedElsewhere{} Moriwaki~\cite{Moriwaki26b}]%
% label removed: prop:thqg-X-heisenberg-sewing-envelope%
\index{Heisenberg!sewing envelope}%
\index{sewing envelope!Heisenberg}%
The sewing envelope of the algebraic Heisenberg VOA
$\cH_\kappa$ with level~$\kappa$ is
\begin{equation}% label removed: eq:thqg-X-heisenberg-sew
\cH_\kappa^{\mathrm{sew}}
\;\cong\;
\operatorname{Sym} A^2(D),
\end{equation}
the symmetric (bosonic) Fock space over the Bergman space of
the disk.  The isomorphism is given by the mode--Bergman
correspondence
\begin{equation}% label removed: eq:thqg-X-mode-bergman-map
\Theta(a_{-n-1}\mathbf{1})
\;=\;
\sqrt{(n+1)\kappa}\; z^n
\;\in\; A^2(D),
\end{equation}
extended multiplicatively by symmetrization:
$\Theta(a_{-n_1-1} \cdots a_{-n_k-1}\mathbf{1})
= \operatorname{Sym}(\Theta(a_{-n_1-1}\mathbf{1})
\otimes \cdots \otimes
\Theta(a_{-n_k-1}\mathbf{1}))$.
\end{proposition}

\begin{proof}
This is the ind-Hilbert identification of
Moriwaki~\cite{Moriwaki26b}, which constructs a conformally
flat 2-disk algebra structure on $\operatorname{Sym} A^2(D)$
and verifies that it is isomorphic to the Hilbert completion
of the Heisenberg Fock space for the sewing-amplitude topology.
The $\sqrt{(n+1)\kappa}$ normalization ensures that
$\|\Theta(a_{-n-1}\mathbf{1})\|^2 = \kappa$, matching the
level of the two-point function
$\langle a(z)\,a(w)\rangle = \kappa/(z-w)$.
\end{proof}

\begin{remark}[Tensor product structure]%
% label removed: rem:thqg-X-tensor-sewing%
The Fock space $\operatorname{Sym} A^2(D)$ has the standard
$n$-particle decomposition
\[
\operatorname{Sym} A^2(D)
\;=\;
\bigoplus_{n=0}^{\infty}
\operatorname{Sym}^n A^2(D),
\]
where $\operatorname{Sym}^0 = \C$ (vacuum sector) and
$\operatorname{Sym}^n$ is the $n$-th symmetric power.
The dimension of the weight-$N$ subspace in
$\operatorname{Sym} A^2(D)$ is the number of partitions
$p(N)$, matching the Heisenberg character
$\chi_{\cH}(q) = \prod_{n=1}^{\infty}(1-q^n)^{-1}
= q^{-1/24}/\dedeta(q)$.
\end{remark}

\subsubsection{Sewing operators on Bergman spaces}

\begin{definition}[Sewing operator]%
% label removed: def:thqg-X-sewing-operator%
\index{sewing operator|textbf}%
\index{Bergman space!sewing operator}%
Let $\Sigma$ be a compact Riemann surface obtained by sewing
two disks $D_1, D_2$ along their boundary circles via the
identification $z_1 z_2 = q$ for $0 < q < 1$.  The
\emph{sewing operator} $\sewop_q \colon A^2(D) \to A^2(D)$ is
defined on the orthonormal basis by
\begin{equation}% label removed: eq:thqg-X-sewing-eigenvalues
\sewop_q\, e_n \;=\; q^{n+1}\, e_n,
\qquad n \geq 0.
\end{equation}
Equivalently, $\sewop_q$ acts by multiplication by
$|z|^2 \cdot q$ on the Bergman kernel:
\begin{equation}% label removed: eq:thqg-X-sewing-kernel
(\sewop_q f)(z)
\;=\;
\int_D \sewop_q(z,w)\,f(w)\,dA(w),
\qquad
\sewop_q(z,w)
\;=\;
\sum_{n=0}^{\infty}
q^{n+1}\,e_n(z)\overline{e_n(w)}.
\end{equation}
\end{definition}

\begin{lemma}[Schatten class properties of $\sewop_q$;
\ClaimStatusProvedHere]%
% label removed: lem:thqg-X-sewing-schatten%
\index{sewing operator!Schatten class}%
For $0 < q < 1$, the sewing operator $\sewop_q$ on $A^2(D)$
satisfies:
\begin{enumerate}[label=\textup{(\roman*)}]
\item \emph{Trace class.}
  $\sewop_q \in \trcl(A^2(D))$ with
  \begin{equation}% label removed: eq:thqg-X-trace-norm
  \|\sewop_q\|_1
  \;=\;
  \sum_{n=0}^{\infty} q^{n+1}
  \;=\;
  \frac{q}{1-q}.
  \end{equation}
\item \emph{Hilbert--Schmidt.}
  $\sewop_q \in \hscl(A^2(D))$ with
  \begin{equation}% label removed: eq:thqg-X-hs-norm
  \|\sewop_q\|_{\hscl}^2
  \;=\;
  \sum_{n=0}^{\infty} q^{2(n+1)}
  \;=\;
  \frac{q^2}{1-q^2}.
  \end{equation}
\item \emph{$p$-Schatten.}
  For all $p \geq 1$,
  $\sewop_q \in \mathcal{L}^p(A^2(D))$ with
  $\|\sewop_q\|_p^p = q^p/(1-q^p)$.
\item \emph{Spectral gap.}
  $\|\sewop_q\| = q$ (operator norm), and the spectral
  radius satisfies $r(\sewop_q) = q < 1$.
\end{enumerate}
\end{lemma}

\begin{proof}
The eigenvalues of $\sewop_q$ are $\lambda_n = q^{n+1}$ for
$n \geq 0$.  Since $0 < q < 1$, the eigenvalues are positive
and summable:
\[
\|\sewop_q\|_p^p
= \sum_{n=0}^{\infty} |\lambda_n|^p
= \sum_{n=0}^{\infty} q^{p(n+1)}
= \frac{q^p}{1-q^p}.
\]
For $p = 1$ (trace norm): $q/(1-q)$.
For $p = 2$ (HS norm squared): $q^2/(1-q^2)$.
The operator norm is $\|\sewop_q\| = \sup_n q^{n+1} = q$.
\end{proof}

\begin{definition}[Genus-$g$ sewing kernel]%
% label removed: def:thqg-X-genus-g-sewing%
\index{sewing operator!genus $g$}%
A genus-$g$ surface $\Sigma_g$ with $n$ punctures admits a
pair-of-pants decomposition into $2g - 2 + n$ pairs of pants,
with $3g - 3 + n$ sewing circles.  Each sewing circle
$C_\alpha$ ($\alpha = 1, \ldots, 3g-3+n$) contributes a sewing
parameter $q_\alpha = e^{-t_\alpha + i\theta_\alpha}$ and
a sewing operator $\sewop_{q_\alpha}$ on $A^2(D)$.
The \emph{genus-$g$ sewing kernel} is the composite operator
\begin{equation}% label removed: eq:thqg-X-genus-g-kernel
\sewker_g
\;=\;
\prod_{\alpha=1}^{3g-3}
\sewop_{q_\alpha}
\;\in\;
\trcl(A^2(D)^{\otimes g}),
\end{equation}
where the product is over the $3g-3$ internal sewing circles
of a genus-$g$ surface without punctures.
More precisely, $\sewker_g$ acts on the tensor product of
$g$ copies of $A^2(D)$ (one for each handle),
with each $\sewop_{q_\alpha}$ acting on the appropriate
tensor factor.
\end{definition}

\begin{remark}[Dependence on pants decomposition]%
% label removed: rem:thqg-X-pants-independence%
\index{pants decomposition!independence}%
The Fredholm determinant $\det(1 - \sewker_g)$ is independent
of the choice of pair-of-pants decomposition, although the
intermediate operator $\sewker_g$ depends on the decomposition.
This independence follows from the sewing axiom of the
chain-level modular-functor package
(Theorem~\ref*{V1-thm:chain-modular-functor}\,(iv)):
different decompositions are related by chain homotopies
that preserve the trace.
\end{remark}

\subsubsection{Composition operators from the prime form}

\begin{definition}[Bergman composition operator]%
% label removed: def:thqg-X-bergman-composition%
\index{composition operator!Bergman}%
Let $\Sigma_g$ be a genus-$g$ Riemann surface with period
matrix $\Omega \in \Siegel_g$ (the Siegel upper half-space),
and let $D_j \subset \Sigma_g$ be small coordinate disks around
marked points $p_j$.  The \emph{Bergman composition operator}
$R \colon A^2(D_{\mathrm{out}}) \to
A^2(D_{\mathrm{in},1}) \otimes A^2(D_{\mathrm{in},2})$
for a pair-of-pants $P$ with one outgoing and two incoming
boundaries is
\begin{equation}% label removed: eq:thqg-X-bergman-composition
R_{n,m}
\;=\;
\frac{1}{(2\pi i)^2}
\oint_{|z|=1} \oint_{|w|=1}
z^{-n-1}\,w^{-m-1}\,
G_P(z,w)\,dz\,dw,
\end{equation}
where $G_P(z,w) = \partial_z \log \primeform(z,w;\Omega)$ is the
chiral Green function on the pair of pants.
\end{definition}

\begin{lemma}[Exponential decay of composition coefficients;
\ClaimStatusProvedHere]%
% label removed: lem:thqg-X-composition-decay%
\index{composition operator!exponential decay}%
The matrix coefficients of the Bergman composition operator satisfy
\begin{equation}% label removed: eq:thqg-X-composition-bound
|R_{n,m}|
\;\leq\;
C\,q^{(n+m)/2}
\end{equation}
for a constant $C$ depending only on the collar length
$t = -\log q$ of the pair of pants.  In particular,
$R$ is Hilbert--Schmidt with
$\|R\|_{\hscl}^2 \leq C^2/(1-q)^2$.
\end{lemma}

\begin{proof}
The Green function $G_P(z,w)$ on the pair of pants is holomorphic
for $|z| < 1$, $|w| < 1$ away from the diagonal $z = w$.
By the collar theorem for hyperbolic surfaces, there exists a
collar of width $t/2$ around each boundary circle, so $G_P$ extends
holomorphically to $|z| < e^{t/2}$, $|w| < e^{t/2}$.

\emph{Contour-moving estimate.}
The matrix coefficient~\eqref{eq:thqg-X-bergman-composition}
is a double contour integral on $|z| = 1$, $|w| = 1$.
Move both contours inward to $|z| = q^{1/2}$, $|w| = q^{1/2}$
(with $q = e^{-t}$, so $q^{1/2} = e^{-t/2} < 1$).
By Cauchy's theorem the integrals are unchanged (the
integrand is holomorphic in the annular region).  On the new
contours:
\begin{itemize}
\item $|z^{-n-1}| = q^{-(n+1)/2}$ and
  $|w^{-m-1}| = q^{-(m+1)/2}$;
\item $|G_P(z,w)| \leq C'$ where
  $C' := \sup_{|z|=q^{1/2},\,|w|=q^{1/2}} |G_P(z,w)|
  < \infty$ (holomorphicity in the collar);
\item each arc has length $2\pi q^{1/2}$.
\end{itemize}
Therefore:
\begin{align*}
|R_{n,m}|
&\leq \frac{1}{(2\pi)^2}
  \cdot (2\pi q^{1/2})^2
  \cdot q^{-(n+1)/2}\,q^{-(m+1)/2}
  \cdot C' \\
&= C'\,q^{1}\,q^{-(n+1)/2}\,q^{-(m+1)/2} \\
&= C'\,q^{(n+m)/2}\,
  \underbrace{q^{1 - (n+1)/2 - (m+1)/2 + (n+m)/2}}_{= q^0 = 1}
  \\
&= C'\,q^{(n+m)/2}.
\end{align*}
(The exponent simplifies:
$1 - \frac{n+1}{2} - \frac{m+1}{2} + \frac{n+m}{2}
= 1 - 1 = 0$.)
Setting $C = C'$ gives the claimed bound.

For the HS norm:
\[
\|R\|_{\hscl}^2
= \sum_{n,m \geq 0} |R_{n,m}|^2
\leq C^2 \sum_{n,m \geq 0} q^{n+m}
= \frac{C^2}{(1-q)^2}.
\qedhere
\]
\end{proof}

\begin{proposition}[Second quantization;
\ClaimStatusProvedHere]%
% label removed: prop:thqg-X-second-quantization%
\index{second quantization!Bergman}%
Let $T \colon A^2(D) \to A^2(D)$ be a trace-class operator
with $\|T\| < 1$.  The \emph{second quantization}
$\secquant(T) \colon \operatorname{Sym} A^2(D) \to
\operatorname{Sym} A^2(D)$ is defined by
\begin{equation}% label removed: eq:thqg-X-second-quant
\secquant(T)|_{\operatorname{Sym}^n A^2(D)}
\;=\;
T^{\otimes n}\big|_{\operatorname{Sym}^n},
\end{equation}
with $\secquant(T)|_{\operatorname{Sym}^0} = \id$.
Then:
\begin{enumerate}[label=\textup{(\roman*)}]
\item $\secquant(T)$ is trace class on $\operatorname{Sym} A^2(D)$
  with
  \begin{equation}% label removed: eq:thqg-X-second-quant-trace
  \operatorname{Tr}_{\operatorname{Sym}}(\secquant(T))
  \;=\;
  \det(1 - T)^{-1}.
  \end{equation}
\item More generally, for the Fredholm determinant:
  \begin{equation}% label removed: eq:thqg-X-fredholm-second-quant
  \det_{\operatorname{Sym}}(1 - \secquant(T))
  \;=\;
  \exp\!\Bigl(-\sum_{k=1}^{\infty}
  \frac{1}{k}\operatorname{Tr}(T^k)\Bigr)
  \;=\;
  \det(1 - T).
  \end{equation}
\end{enumerate}
\end{proposition}

\begin{proof}
Part~(i): The trace on $\operatorname{Sym}^n A^2(D)$ is computed
using an orthonormal basis of symmetric tensors.
For eigenvalues $\lambda_j$ of~$T$:
\[
\operatorname{Tr}_{\operatorname{Sym}^n}(T^{\otimes n})
= \sum_{j_1 \leq \cdots \leq j_n}
\lambda_{j_1} \cdots \lambda_{j_n}
= h_n(\lambda_1, \lambda_2, \ldots),
\]
where $h_n$ is the complete symmetric polynomial.
Summing over $n \geq 0$:
\[
\sum_{n=0}^{\infty}
\operatorname{Tr}_{\operatorname{Sym}^n}(T^{\otimes n})
= \sum_{n=0}^{\infty} h_n(\lambda)
= \prod_{j} \frac{1}{1 - \lambda_j}
= \det(1 - T)^{-1}.
\]
The series converges absolutely because $T$ is trace class
and $\|T\| < 1$, so $\sum_j |\lambda_j| < \infty$ and
$|\lambda_j| < 1$ for all~$j$.

Part~(ii): $\log\det(1-T) = \sum_j \log(1-\lambda_j)
= -\sum_{k=1}^{\infty} \frac{1}{k}\sum_j \lambda_j^k
= -\sum_{k=1}^{\infty} \frac{1}{k}\operatorname{Tr}(T^k)$.
\end{proof}

% =====================================================================
%
%  SECTION 2: THE HEISENBERG SEWING THEOREM
%
% =====================================================================

\subsection{The Heisenberg sewing theorem: full development}
% label removed: subsec:thqg-X-heisenberg-sewing

\subsubsection{Statement and architecture}

The Heisenberg sewing theorem
(Theorem~\ref{thm:heisenberg-sewing})
has four clauses.  We now develop the complete proofs, starting
from the mode--Bergman correspondence and building to the
Fredholm determinant formula for genus-$g$ partition functions.

\begin{theorem}[Heisenberg sewing theorem: full development;
\ClaimStatusProvedHere]%
% label removed: thm:thqg-X-heisenberg-sewing-full%
\index{Heisenberg!sewing theorem!full development}%
Let $\cH_\kappa$ be the Heisenberg algebra with level $\kappa > 0$,
equipped with the mode--Bergman correspondence~$\Theta$
\textup{(}Remark~\textup{\ref{rem:heisenberg-mode-bergman}}\textup{)}.
Then:
\begin{enumerate}[label=\textup{(\Roman*)}]
\item \textbf{Sewing envelope identification.}
  The sewing envelope of $\cH_\kappa$ is
  $\cH_\kappa^{\mathrm{sew}} \cong
  \operatorname{Sym} A^2(D)$.
\item \textbf{Bar differential closure.}
  The algebraic bar differential extends by closure to the
  Gaussian collision operator on Bergman tensors:
  \begin{equation}% label removed: eq:thqg-X-bar-closure
  d_{\mathrm{bar}}^{\mathrm{an}}
  \;=\;
  \overline{d_{\mathrm{bar}}^{\mathrm{alg}}}
  ^{\,\graphnorm}
  \;=\;
  \sum_{i < j}
  \kappa\;\operatorname{Res}_{z_i = z_j}
  \frac{1}{(z_i - z_j)}\,
  \partial_{z_i z_j}.
  \end{equation}
\item \textbf{Pair-of-pants sewing = second quantization.}
  The pair-of-pants amplitude on
  $\operatorname{Sym} A^2(D)$ is the second quantization
  $\secquant(R)$ of the one-particle Bergman restriction
  operator
  $R \colon A^2(D_{\mathrm{out}}) \to
  A^2(D_{\mathrm{in},1}) \otimes A^2(D_{\mathrm{in},2})$.
\item \textbf{Fredholm determinant formula.}
  Every genus-$g$ closed amplitude is a Fredholm determinant:
  the genus-$g$ partition function of $\cH_\kappa$ is
  \begin{equation}% label removed: eq:thqg-X-fredholm-formula
  Z_g(\cH_\kappa;\,\Omega)
  \;=\;
  \bigl(\det\operatorname{Im}\Omega\bigr)^{-\kappa/2}
  \cdot
  \bigl(\zetareg\Delta_{\Sigma_g}\bigr)^{-\kappa/2},
  \end{equation}
  where $\Omega$ is the period matrix of $\Sigma_g$ and
  $\zetareg\Delta$ is the zeta-regularized determinant of the
  Laplacian.
\end{enumerate}
\end{theorem}

The proof occupies the remainder of this subsection.

\subsubsection{Clause I: sewing envelope identification}

\begin{proof}[Proof of Clause~I]
The mode--Bergman map $\Theta$ of
Remark~\ref{rem:heisenberg-mode-bergman} defines a dense
algebraic embedding
$\Theta \colon \Fock_\kappa^{\mathrm{fin}} \hookrightarrow
\operatorname{Sym} A^2(D)$
from the algebraic Fock space (finite linear combinations of
mode monomials) into the symmetric algebra of the Bergman space.

\emph{Step 1: algebraic amplitudes.}
On the algebraic Fock core, every amplitude
$A_\Sigma^{\mathrm{alg}}$ is a finite sum of Wick contractions.
By Wick's theorem for the free boson, the amplitude of
$a_{-n_1-1}\cdots a_{-n_k-1}\mathbf{1}$ on a surface~$\Sigma$
with collar coordinates is a polynomial in the mode
coefficients of the propagator $G_\Sigma$, which by the
mode--Bergman dictionary become matrix coefficients of the
Bergman composition operator $R$.

\emph{Step 2: continuity for the sewing topology.}
The sewing seminorms $p_{\Sigma,\xi,\eta}$ on $\Fock_\kappa^{\mathrm{fin}}$
are controlled by the Bergman norms: for
$v = \Theta(a_{-n_1-1}\cdots a_{-n_k-1}\mathbf{1})
\in \operatorname{Sym}^k A^2(D)$,
\begin{equation}% label removed: eq:thqg-X-sewing-seminorm-bound
p_{\Sigma,\xi,\eta}(v)
\;\leq\;
\|\xi\| \cdot \|\eta\| \cdot \|v\|_{\operatorname{Sym}^k}
\cdot \prod_\alpha \|\sewop_{q_\alpha}\|,
\end{equation}
where the product is over sewing circles of~$\Sigma$.
Since each $\sewop_{q_\alpha}$ is bounded with
$\|\sewop_{q_\alpha}\| = q_\alpha < 1$, the sewing topology is
weaker than the $\operatorname{Sym} A^2(D)$ topology.

\emph{Step 3: completeness.}
Conversely, any Cauchy sequence in the sewing topology is Cauchy
in $\operatorname{Sym} A^2(D)$ (because the latter norm
controls each sewing seminorm).  Hence the completions coincide:
$\cH_\kappa^{\mathrm{sew}} \cong \operatorname{Sym} A^2(D)$.
The full argument is the identification of Moriwaki~\cite{Moriwaki26b},
which constructs the conformally flat 2-disk algebra structure
on $\operatorname{Sym} A^2(D)$ and verifies the universal
property of the sewing envelope.
\end{proof}

\subsubsection{Clause II: bar differential closure}

\begin{proof}[Proof of Clause~II]
The algebraic bar differential on $\barB_n(\cH_\kappa)$ is
\[
d_{\mathrm{bar}}^{\mathrm{alg}}
= \sum_{i < j}
\operatorname{Res}_{z_i = z_j}
\circ\, Y_{ij},
\]
where $Y_{ij}$ is the chiral multiplication
$Y(a_{-n-1}\mathbf{1}, z_i - z_j)
= \sum_m a_m (z_i - z_j)^{-m-1}$.
For the Heisenberg algebra, the OPE is
$a(z_i)\,a(z_j) \sim \kappa/(z_i - z_j)$,
so $Y_{ij}$ acts on the one-particle Bergman space as
multiplication by $\kappa/(z_i - z_j)$ followed by the residue
at $z_i = z_j$.

\emph{Step 1: densely defined operator.}
On $\operatorname{Sym}^n A^2(D)$, the operator
$d_{\mathrm{bar}}^{\mathrm{alg}}$ is densely defined on the
algebraic Fock core.  Its action on monomials is
\begin{align}
d_{\mathrm{bar}}(z^{n_1} \cdots z^{n_k})
&= \kappa \sum_{i < j}
\operatorname{Res}_{z_i = z_j}
\frac{z_1^{n_1} \cdots z_k^{n_k}}{z_i - z_j}
\notag\\
&= \kappa \sum_{i < j}
\sum_{p=0}^{n_j-1}
z_j^{n_i + n_j - 1 - p}\, z_j^p
\cdot \prod_{\ell \neq i,j} z_\ell^{n_\ell},
% label removed: eq:thqg-X-bar-monomial-action
\end{align}
where the residue replaces $z_i$ by $z_j$ with the standard
contraction.

\emph{Step 2: closability.}
The operator $d_{\mathrm{bar}}^{\mathrm{alg}}$ is closable
on $\operatorname{Sym}^n A^2(D)$.  This follows from the
fact that the formal adjoint $(d_{\mathrm{bar}}^{\mathrm{alg}})^*$
is also densely defined (it is the creation operator that reverses
the contraction), so the closure exists by von Neumann's theorem.

\emph{Step 3: identification of the closure.}
The closure is the \emph{Gaussian collision operator}: the sum
of pairwise collision residue operators weighted by the level
$\kappa$.  Explicitly, it acts on
$f_1 \otimes \cdots \otimes f_n \in \operatorname{Sym}^n A^2(D)$
by
\[
\overline{d}_{\mathrm{bar}}(f_1 \otimes \cdots \otimes f_n)
= \kappa \sum_{i < j}
\operatorname{Res}_{z_i = z_j}
\frac{f_i(z_i)\,f_j(z_j)}{z_i - z_j}
\cdot \prod_{\ell \neq i,j} f_\ell(z_\ell),
\]
where the residue is computed by the Bergman projection
of $f_i(z)\,f_j(z) / (z_i - z_j)$ restricted to the diagonal.

To see that $\overline{d}_{\mathrm{bar}}$ is closed in
the Bergman norm topology, consider the $(i,j)$ summand:
the collision operator $C_{ij}(f_i \otimes f_j)(z)
= \operatorname{Res}_{w=z}[f_i(w)f_j(z)/(w-z)]$.
On the orthonormal basis $e_n(z) = \sqrt{(n+1)/\pi}\,z^n$:
\[
C_{ij}(e_a \otimes e_b)(z)
= \operatorname{Res}_{w=z}\!
\frac{\sqrt{(a+1)(b+1)}}{\pi}
\frac{w^a z^b}{w - z}
= \frac{\sqrt{(a+1)(b+1)}}{\pi}
\sum_{p=0}^{a-1} z^{a+b-1-p}\cdot z^p
\]
by geometric series expansion of $w^a/(w-z)$.  The
Bergman norm of the output satisfies
$\|C_{ij}(e_a \otimes e_b)\|^2
\leq a \cdot (a+1)(b+1)/\pi^2 \cdot (\text{sum of norms})$,
which is bounded by $a(a+1)(b+1)/\pi \leq (a+1)^2(b+1)$.
The graph norm
$\|x\|_{\mathrm{graph}}^2
:= \|x\|^2 + \|\overline{d}_{\mathrm{bar}} x\|^2$
is therefore controlled by the Bergman norm of the input
(with polynomial growth in the mode numbers), establishing
$\overline{d}_{\mathrm{bar}}$ as a closed operator.  By
the uniqueness of closure for closable operators (von Neumann),
$\overline{d}_{\mathrm{bar}}$ is the unique closed extension
of $d_{\mathrm{bar}}^{\mathrm{alg}}$.
\end{proof}

\subsubsection{Clause III: pair-of-pants = second quantization}

\begin{proof}[Proof of Clause~III]
Consider a pair of pants $P$ with one outgoing boundary $\partial_{\mathrm{out}}$
and two incoming boundaries $\partial_{\mathrm{in},1}$, $\partial_{\mathrm{in},2}$.
Each boundary is parametrized by a circle, and the sewing
parameters $(q_1, q_2)$ are the collar parameters.

\emph{Step 1: one-particle amplitude.}
On the one-particle space $A^2(D)$, the pair-of-pants amplitude
is the Bergman composition operator
$R \colon A^2(D_{\mathrm{out}}) \to
A^2(D_{\mathrm{in},1}) \otimes A^2(D_{\mathrm{in},2})$
of Definition~\ref{def:thqg-X-bergman-composition}.
Its matrix coefficients $R_{n,m}$ are the contour integrals
of the pair-of-pants Green function $G_P(z,w)$.

\emph{Step 2: Wick's theorem.}
On the full Fock space $\operatorname{Sym} A^2(D)$, the
Heisenberg amplitudes are determined by Wick's theorem:
every $n$-point function is a sum over pairings of two-point
functions.  In the Bergman picture, the $n$-particle amplitude
on $P$ is
\begin{equation}% label removed: eq:thqg-X-wick-bergman
A_P^{(n)}(f_1 \otimes \cdots \otimes f_n)
= \sum_{\sigma \in S_n}
\prod_{j=1}^{n}
\langle f_j, R\,g_{\sigma(j)}\rangle,
\end{equation}
where $g_{\sigma(j)}$ are the incoming modes.  This is the formula for the second quantization:
\[
A_P^{(n)}
= \secquant(R)\big|_{\operatorname{Sym}^n}.
\]

\emph{Step 3: total amplitude.}
Summing over particle number gives the full pair-of-pants
amplitude:
\begin{equation}% label removed: eq:thqg-X-pants-second-quant
A_P
= \bigoplus_{n=0}^{\infty}
A_P^{(n)}
= \secquant(R)
\colon \operatorname{Sym} A^2(D_{\mathrm{out}})
\to
\operatorname{Sym}(A^2(D_{\mathrm{in},1})
\oplus A^2(D_{\mathrm{in},2})).
\end{equation}
The convergence of this series follows from the exponential
decay $|R_{n,m}| \leq C q^{(n+m)/2}$
(Lemma~\ref{lem:thqg-X-composition-decay}).
\end{proof}

\subsubsection{Clause IV: Fredholm determinant}

\begin{proof}[Proof of Clause~IV]
A genus-$g$ surface $\Sigma_g$ without punctures is obtained by
sewing together $2g-2$ pairs of pants along $3g-3$ sewing circles.
Each sewing circle contributes a trace over the one-particle
Hilbert space, and by Wick's theorem the total amplitude factors
through the one-particle operators.

\emph{Step 1: genus-$g$ amplitude as iterated trace.}
The genus-$g$ partition function is
\begin{equation}% label removed: eq:thqg-X-genus-g-trace
Z_g(\cH_\kappa; \Omega)
= \operatorname{Tr}_{\operatorname{Sym} A^2}
\Bigl(\prod_{\alpha=1}^{3g-3}
\secquant(\sewop_{q_\alpha})^{\kappa}\Bigr),
\end{equation}
where the product is over all sewing circles, $\sewop_{q_\alpha}$
is the one-particle sewing operator, and the exponent $\kappa$
reflects the level (one Fock space per unit of central charge).

\emph{Step 2: trace factorization.}
By Proposition~\ref{prop:thqg-X-second-quantization}\,(i):
\[
\operatorname{Tr}_{\operatorname{Sym}}(\secquant(T))
= \det(1 - T)^{-1}
\]
for any trace-class $T$ with $\|T\| < 1$.  Applying this to each
sewing circle:
\begin{equation}% label removed: eq:thqg-X-genus-g-det
Z_g(\cH_\kappa; \Omega)
= \prod_{\alpha=1}^{3g-3}
\det\bigl(1 - \sewop_{q_\alpha}\bigr)^{-\kappa}.
\end{equation}

\emph{Step 3: from pants parameters to period matrix.}
The sewing parameters $\{q_\alpha\}$ are the Fenchel--Nielsen
coordinates of $\Sigma_g$; they determine and are determined by
the period matrix $\Omega \in \Siegel_g$.  The Polyakov--Alvarez
formula \cite{BK86} relates the product of one-particle
determinants to the period matrix and the zeta-regularized
determinant of the Laplacian:
\begin{equation}% label removed: eq:thqg-X-polyakov
\prod_{\alpha=1}^{3g-3}
\det(1 - \sewop_{q_\alpha})^{-1}
= (\det\operatorname{Im}\Omega)^{-1/2}
\cdot (\zetareg\Delta_{\Sigma_g})^{-1/2}.
\end{equation}
This is the holomorphic factorization of the string measure
of Belavin--Knizhnik~\cite{BK86}: the modular-invariant
expression factors into a product of a holomorphic and an
anti-holomorphic part, and the holomorphic part is
$\det(1 - \sewker_g)^{-1}$ up to the period-matrix prefactor.

\emph{Step 4: level $\kappa$.}
Raising to the power $\kappa$ (i.e., tensoring $\kappa$ copies
of the one-particle Bergman space, or equivalently working at
level $\kappa$):
\begin{equation}% label removed: eq:thqg-X-final-fredholm
Z_g(\cH_\kappa; \Omega)
= (\det\operatorname{Im}\Omega)^{-\kappa/2}
\cdot (\zetareg\Delta_{\Sigma_g})^{-\kappa/2}.
\end{equation}
This completes the proof of Clause~IV and of the full theorem.
\end{proof}

% =====================================================================
%
%  SECTION 3: ONE-PARTICLE REDUCTION
%
% =====================================================================

\subsection{One-particle reduction: trace-class operators
and the Dedekind eta}
% label removed: subsec:thqg-X-one-particle

The one-particle reduction principle reduces the computation
of partition functions on the bosonic Fock space to Fredholm
determinants on the one-particle Bergman space.  This subsection
carries out the explicit computations at genus $1$, $2$, and $3$.

\subsubsection{Genus 1: the torus and the Dedekind eta}

\begin{theorem}[Genus-1 Fredholm determinant;
\ClaimStatusProvedHere]%
% label removed: thm:thqg-X-genus-1-fredholm%
\index{Fredholm determinant!genus 1}%
\index{Dedekind eta function!from Fredholm determinant}%
The genus-$1$ partition function of $\cH_\kappa$ is
\begin{equation}% label removed: eq:thqg-X-genus-1-det
Z_1(\cH_\kappa;\,\tau)
\;=\;
(\operatorname{Im}\tau)^{-\kappa/2}\,
\bigl|\dedeta(\tau)\bigr|^{-2\kappa},
\end{equation}
where the holomorphic part is
\begin{equation}% label removed: eq:thqg-X-genus-1-hol
Z_1^{\mathrm{hol}}(\cH_\kappa;\,\tau)
\;=\;
\dedeta(\tau)^{-\kappa}
\;=\;
q^{-\kappa/24}\,
\prod_{n=1}^{\infty}
(1 - q^n)^{-\kappa}.
\end{equation}
\end{theorem}

\begin{proof}
A genus-$1$ surface $E_\tau = \C/(\Z + \tau\Z)$ is obtained
by sewing a single cylinder $\C^* \supset
\{q < |z| < 1\}$ with $q = e^{2\pi i \tau}$ ($\operatorname{Im}\tau > 0$).
The sewing circle is a single boundary component, and the
sewing operator is $\sewop_q$ on $A^2(D)$ with eigenvalues
$q^{n+1}$.

\emph{Step 1: Fock space trace.}
The partition function is the Fock-space trace:
\begin{align}
Z_1(\cH_1;\,\tau)
&= \operatorname{Tr}_{\operatorname{Sym} A^2(D)}
\bigl(\secquant(\sewop_q)\bigr)
\notag\\
&= \det(1 - \sewop_q)^{-1}
\quad\text{(Proposition~\ref{prop:thqg-X-second-quantization}\,(i))}
\notag\\
&= \prod_{n=0}^{\infty}
(1 - q^{n+1})^{-1}
\notag\\
&= \prod_{n=1}^{\infty}
(1 - q^n)^{-1}.
% label removed: eq:thqg-X-genus-1-product
\end{align}

\emph{Step 2: Dedekind eta identification.}
The Dedekind eta function is
$\dedeta(\tau) = q^{1/24}\prod_{n=1}^{\infty}(1-q^n)$
(Definition~\ref{def:eta-function}).  Therefore:
\begin{equation}% label removed: eq:thqg-X-eta-identification
\prod_{n=1}^{\infty}(1-q^n)^{-1}
= q^{1/24} / \dedeta(\tau)
= q^{1/24}\,\dedeta(\tau)^{-1}.
\end{equation}
Including the vacuum energy $q^{-c/24} = q^{-1/24}$ for
$c = 1$ (one free boson with central charge $1$ per unit of
$\kappa$) and the non-holomorphic prefactor
$(\operatorname{Im}\tau)^{-1/2}$ from the zero-mode
integral:
\[
Z_1(\cH_1;\,\tau)
= (\operatorname{Im}\tau)^{-1/2}\,
q^{-1/24}\,\bar{q}^{-1/24}\,
\prod_{n=1}^{\infty}
\frac{1}{|1-q^n|^2}
= (\operatorname{Im}\tau)^{-1/2}\,
|\dedeta(\tau)|^{-2}.
\]

\emph{Step 3: level $\kappa$.}
At level $\kappa$, the Fock space is the $\kappa$-fold tensor
product of the $\kappa = 1$ Fock space (or equivalently, the
sewing operator eigenvalues are raised to the power $\kappa$):
\[
Z_1(\cH_\kappa;\,\tau)
= (\operatorname{Im}\tau)^{-\kappa/2}\,
|\dedeta(\tau)|^{-2\kappa}.
\qedhere
\]
\end{proof}

\begin{remark}[Modular transformation]%
% label removed: rem:thqg-X-genus-1-modular%
Under $\tau \mapsto -1/\tau$:
$\dedeta(-1/\tau) = \sqrt{-i\tau}\,\dedeta(\tau)$, so
\[
Z_1(\cH_\kappa;\,-1/\tau)
= |\tau|^{-\kappa}\,
(\operatorname{Im}\tau/|\tau|^2)^{-\kappa/2}\,
|\dedeta(\tau)|^{-2\kappa}
= Z_1(\cH_\kappa;\,\tau).
\]
The $S$-transformation invariance is automatic from the
Fredholm determinant structure: $\det(1 - \sewop_q)$ depends
only on the eigenvalues, which are determined by the conformal
structure of $E_\tau$, not the choice of fundamental domain.
\end{remark}

\begin{computation}[Genus-1 Fredholm determinant: explicit series;
\ClaimStatusProvedHere]%
% label removed: comp:thqg-X-genus-1-series%
\index{Fredholm determinant!genus 1!explicit series}%
The Fredholm determinant at genus~$1$ has the $q$-expansion:
\begin{align}
\det(1 - \sewop_q)^{-1}
&= \prod_{n=1}^{\infty} (1 - q^n)^{-1}
\notag\\
&= 1 + q + 2q^2 + 3q^3 + 5q^4 + 7q^5 + 11q^6 + 15q^7
+ 22q^8 + \cdots
\notag\\
&= \sum_{n=0}^{\infty} p(n)\,q^n,
% label removed: eq:thqg-X-partition-numbers
\end{align}
where $p(n)$ is the number of integer partitions of~$n$.
The asymptotic behavior is given by the Hardy--Ramanujan formula:
\begin{equation}% label removed: eq:thqg-X-hardy-ramanujan
p(n) \;\sim\;
\frac{1}{4n\sqrt{3}}\,
\exp\!\Bigl(\pi\sqrt{2n/3}\Bigr)
\quad\text{as } n \to \infty.
\end{equation}
The subexponential growth $\log p(n) \sim \pi\sqrt{2n/3}$
is precisely the condition needed for HS-sewing
(Theorem~\ref{thm:general-hs-sewing}).
\end{computation}

\begin{remark}[Trace-class vs.\ Fredholm: comparison of
convergence]%
% label removed: rem:thqg-X-trace-vs-fredholm%
The trace $\sum_n q^{n+1} = q/(1-q)$ of $\sewop_q$ converges
for $|q| < 1$, as does every $p$-Schatten norm.
The Fredholm determinant
$\det(1 - \sewop_q)^{-1}
= \prod_n (1 - q^n)^{-1}$
converges absolutely for $|q| < 1$ (the Euler function).
The radius of convergence in~$q$ is $|q| = 1$, where the
product diverges because $\sum_n q^n$ diverges.
On the modular side, $|q| = 1$ corresponds to
$\operatorname{Im}\tau = 0$, the cusp.
\end{remark}

\subsubsection{Genus 2: two handles and the Siegel modular form}

\begin{theorem}[Genus-2 Fredholm determinant;
\ClaimStatusProvedHere]%
% label removed: thm:thqg-X-genus-2-fredholm%
\index{Fredholm determinant!genus 2}%
The genus-$2$ partition function of $\cH_\kappa$ on a
surface $\Sigma_2$ with period matrix
$\Omega = \bigl(\begin{smallmatrix}\tau_1 & z \\ z & \tau_2
\end{smallmatrix}\bigr)
\in \Siegel_2$ is
\begin{equation}% label removed: eq:thqg-X-genus-2-det
Z_2(\cH_\kappa;\,\Omega)
\;=\;
\bigl(\det\operatorname{Im}\Omega\bigr)^{-\kappa/2}
\cdot
\bigl(\zetareg\Delta_{\Sigma_2}\bigr)^{-\kappa/2}.
\end{equation}
The holomorphic part is the Fredholm determinant of the
genus-$2$ sewing kernel $\sewker_2$ acting on $A^2(D)^{\oplus 2}$:
\begin{equation}% label removed: eq:thqg-X-genus-2-hol
Z_2^{\mathrm{hol}}(\cH_1;\,\Omega)
\;=\;
\det\bigl(1 - \sewker_2\bigr)^{-1},
\end{equation}
where $\sewker_2$ is the $2 \times 2$ matrix operator on
$A^2(D)^{\oplus 2}$ with entries determined by the sewing data
of the pair-of-pants decomposition.
\end{theorem}

\begin{proof}
A genus-$2$ surface admits a pair-of-pants decomposition into
$2$ pairs of pants ($2g-2 = 2$) with $3$ sewing circles
($3g-3 = 3$).

\emph{Step 1: pants decomposition.}
Fix a pair-of-pants decomposition: $\Sigma_2 = P_1 \cup_{\partial} P_2$,
where $P_1$ and $P_2$ are each pairs of pants (three-holed spheres).
The three sewing circles have parameters $q_1, q_2, q_3$, where
$q_1$ and $q_2$ correspond to the two handles (B-cycles) and
$q_3$ corresponds to the internal seam.

\emph{Step 2: one-particle operator.}
On the one-particle Bergman space, the sewing produces a
composition of three operators.  Write the one-particle sewing
kernel as a $2 \times 2$ block operator on
$A^2(D)^{\oplus 2}$ (one block per handle):
\begin{equation}% label removed: eq:thqg-X-genus-2-block
\sewker_2
=
\begin{pmatrix}
\sewop_{q_1} & R_{12} \\
R_{21} & \sewop_{q_2}
\end{pmatrix},
\end{equation}
where $\sewop_{q_j}$ are the handle sewing operators and
$R_{12}, R_{21}$ are the off-diagonal composition operators
encoding the coupling between the two handles through the internal
seam.  The off-diagonal entries satisfy
$\|R_{12}\|_{\hscl}, \|R_{21}\|_{\hscl} \leq C\,q_3^{1/2}$
by Lemma~\ref{lem:thqg-X-composition-decay}.

\emph{Step 3: Fredholm determinant.}
Since all diagonal and off-diagonal entries are trace class
(the diagonal operators $\sewop_{q_j}$ are trace class by
Lemma~\ref{lem:thqg-X-sewing-schatten}, and the off-diagonal
operators are Hilbert--Schmidt with $\hscl \cdot \hscl \subset \trcl$),
the full operator $\sewker_2$ is trace class on
$A^2(D)^{\oplus 2}$.  The Fredholm determinant is:
\begin{align}
\det(1 - \sewker_2)
&= \det\!\begin{pmatrix}
1 - \sewop_{q_1} & -R_{12} \\
-R_{21} & 1 - \sewop_{q_2}
\end{pmatrix}
\notag\\
&= \det(1 - \sewop_{q_1})\,
\det\bigl(1 - \sewop_{q_2}
- R_{21}(1-\sewop_{q_1})^{-1}R_{12}\bigr).
% label removed: eq:thqg-X-genus-2-schur
\end{align}
The Schur complement
$S = \sewop_{q_2} + R_{21}(1-\sewop_{q_1})^{-1}R_{12}$
captures the interaction between the two handles.

\emph{Step 4: degeneration limits.}
In the separating degeneration $q_3 \to 0$
(pinching the internal seam):
$R_{12}, R_{21} \to 0$, and
\[
\det(1 - \sewker_2) \to
\det(1 - \sewop_{q_1}) \cdot \det(1 - \sewop_{q_2})
= \dedeta(\tau_1)^2 \cdot \dedeta(\tau_2)^2,
\]
recovering the product of two genus-$1$ partition functions.
In the non-separating degeneration $q_1 \to 0$
(collapsing one handle):
$\sewop_{q_1} \to 0$ and $\det(1 - \sewker_2)$ reduces to
$\det(1 - \sewop_{q_2} - R_{21}R_{12})$, a genus-$1$
Fredholm determinant with modified sewing operator.

\emph{Step 5: Polyakov formula.}
Applying the Polyakov--Alvarez formula
\eqref{eq:thqg-X-polyakov} at genus~$2$ gives
$\det(1 - \sewker_2)^{-1}
= (\det\operatorname{Im}\Omega)^{-1/2}
\cdot (\zetareg\Delta_{\Sigma_2})^{-1/2}$.
Raising to the power $\kappa$ yields the stated formula.
\end{proof}

\begin{computation}[Genus-2 Fredholm determinant: leading terms;
\ClaimStatusProvedHere]%
% label removed: comp:thqg-X-genus-2-series%
\index{Fredholm determinant!genus 2!leading terms}%
Write $q_1 = e^{2\pi i \tau_1}$, $q_2 = e^{2\pi i \tau_2}$,
$r = e^{2\pi i z}$ (the off-diagonal period).
In the product expansion:
\begin{align}
\det(1 - \sewker_2)^{-1}
&= \prod_{n=1}^{\infty}(1-q_1^n)^{-1}
\cdot \prod_{n=1}^{\infty}(1-q_2^n)^{-1}
\cdot \bigl(1 + O(r)\bigr)
\notag\\
&= \frac{1}{\dedeta(\tau_1)\,\dedeta(\tau_2)}
\cdot \Bigl(1 + \sum_{n,m \geq 1}
c_{n,m}\,r^n\,\bar{r}^m + \cdots\Bigr).
% label removed: eq:thqg-X-genus-2-expansion
\end{align}
The leading correction $c_{1,1}$ arises from the one-particle
Schur complement and equals
$\sum_{k \geq 0} q_1^{k+1}\,q_2^{k+1}/(1-q_1^{k+1})(1-q_2^{k+1})$.
For $\tau_1 = \tau_2 = i$ (square torus) and $z = i/2$:
$|r| = e^{-\pi}$ and $|\dedeta(i)|^{-2} = \Gamma(1/4)^2/(2\pi^{3/2})$.
\end{computation}

\begin{remark}[Arithmetic significance of genus-$2$
non-diagonality]%
\label{rem:genus2-arithmetic-significance}%
\index{sewing kernel!genus-2 arithmetic significance}%
The off-diagonal blocks $R_{12}, R_{21}$ in the genus-$2$
sewing kernel $K_2$ have a precise arithmetic consequence.
At genus~$1$, the sewing operator is Fock-diagonal, which
forces the two-variable $L$-object $L_\cA(s,u)$ to collapse
to a single-variable Dirichlet series $S_\cA(s{+}u{-}1)$
(Vol~I, sewing--Hecke reciprocity theorem);
this collapse is the structural reason that genus-$1$ MC
data cannot access zero-location information
(Vol~I, structural separation theorem).
At genus~$2$, the off-diagonal coupling $R_{ij} \neq 0$
(for generic period matrix) breaks this diagonality.
The Schur complement
$S = K_{q_2} + R_{21}(1{-}K_{q_1})^{-1}R_{12}$
encodes inter-handle interaction that is invisible to the
genus-$1$ sewing operator; the resulting Fredholm
determinant depends nontrivially on the off-diagonal
period~$z$, and the Rankin--Selberg unfolding on
$\mathrm{Sp}_4(\mathbb{Z}) \backslash \mathfrak{H}_2$
produces genuinely multi-variable $L$-functions.
The genus-$2$ arithmetic programme
(Vol~I, arithmetic shadows chapter,
genus-$2$ arithmetic frontier)
exploits this non-diagonality to escape the genus-$1$
structural barrier.
\end{remark}

\begin{remark}[The four Dirichlet series and the arity--genus orthogonality]%
\label{rem:four-dirichlet-arity-genus}%
\index{Dirichlet series!four-object hierarchy}%
\index{arity--genus orthogonality|textbf}
The genus-$1$ structural barrier of the preceding remark
is a manifestation of a deeper phenomenon: the arithmetic
programme of Volume~I involves four distinct Dirichlet series that
live on orthogonal axes of the bigraded shadow algebra
$\cA^{\mathrm{sh}}_{r,g}$.

\begin{enumerate}[label=\textup{(\alph*)}]
\item The \emph{shadow-metric Epstein zeta}
  $\varepsilon_{Q_L}(s)$ and the \emph{shadow Dirichlet series}
  $\zeta_\cA(s) = \sum_{r \geq 2} S_r\,r^{-s}$
  live on the \textbf{arity axis}~$(r, 0)$: they are
  genus-$0$ invariants determined by the OPE data
  $(\kappa, \alpha, S_4)$.

\item The \emph{constrained Epstein zeta}
  $\varepsilon^c_s = \sum_\Delta d(\Delta)\,(2\Delta)^{-s}$
  (Benjamin--Chang~\cite{BenjaminChang22})
  lives on the \textbf{genus axis}~$(0, 1)$: it is the Eisenstein
  spectral coefficient of the descendant-stripped partition
  function $\hat Z^c = y^{c/2}|\eta|^{2c}\,Z$ on
  $\mathrm{SL}(2,\bZ)\backslash\mathfrak{H}$.

\item The \emph{categorical zeta}
  $\zeta^{\mathrm{DK}}_\fg(s) = \sum_\lambda
  \dim(V_\lambda)^{-s}$ is a purely Lie-theoretic invariant
  ($= \zeta(s)$ for $\fsl_2$, Vol~I,
  Proposition~\ref*{prop:categorical-zeta-sl2-riemann}).
\end{enumerate}

The nontrivial Riemann zeta zeros~$\rho_n$ appear through
object~(b), not~(a): the scattering factor
$\zeta(2s)/\zeta(2s{-}1)$ in the Eisenstein spectral
decomposition on~$\cM_{1,1}$ has poles at
$s = (1 + \rho_n)/2$.  This is a theorem about the spectral
theory of the \emph{moduli space}, not about the bar complex.
The ``negative results''
(shadow zeta $\notin$ Selberg class, no Euler product)
concern object~(a), which lives on the arity axis and has
no access to the Riemann zeta zeros.

The genus-$2$ non-diagonality of the preceding remark
provides the first mechanism for the arity and genus axes to
\emph{interact}: the off-diagonal sewing blocks~$R_{ij}$
couple the handle structure (genus axis) to the inter-handle
OPE data (arity axis), breaking the genus-$1$ factorisation
barrier.  The Dirichlet-sewing lift
$S_\cA(u) = \zeta(u{+}1)\sum_i (\zeta(u) - H_{w_i-1}(u))$
(Vol~I) is a fifth related object, built from the weight
multiset~$W(\cA)$, that contains $\zeta$-factors
($S_\cH(u) = \zeta(u)\zeta(u{+}1)$ for Heisenberg)
but is independent of the shadow obstruction tower.
See Volume~I, \S\ref*{subsec:four-dirichlet-series}
for the full disambiguation.
\end{remark}

\subsubsection{Genus 3: the three-handle Fredholm determinant}

\begin{theorem}[Genus-3 Fredholm determinant;
\ClaimStatusProvedHere]%
% label removed: thm:thqg-X-genus-3-fredholm%
\index{Fredholm determinant!genus 3}%
The genus-$3$ partition function of $\cH_\kappa$ on
$\Sigma_3$ with period matrix $\Omega \in \Siegel_3$ is
\begin{equation}% label removed: eq:thqg-X-genus-3-det
Z_3(\cH_\kappa;\,\Omega)
\;=\;
\bigl(\det\operatorname{Im}\Omega\bigr)^{-\kappa/2}
\cdot
\bigl(\zetareg\Delta_{\Sigma_3}\bigr)^{-\kappa/2}.
\end{equation}
The one-particle sewing kernel $\sewker_3$ acts on
$A^2(D)^{\oplus 3}$ as a $3 \times 3$ operator:
\begin{equation}% label removed: eq:thqg-X-genus-3-block
\sewker_3
=
\begin{pmatrix}
\sewop_{q_1} & R_{12} & R_{13} \\
R_{21} & \sewop_{q_2} & R_{23} \\
R_{31} & R_{32} & \sewop_{q_3}
\end{pmatrix},
\end{equation}
where the diagonal entries $\sewop_{q_j}$ correspond to the
three handles and the off-diagonal entries $R_{ij}$ encode the
pair-of-pants composition operators for the internal seams.
\end{theorem}

\begin{proof}
A genus-$3$ surface admits a pair-of-pants decomposition into
$4$ pairs of pants ($2g - 2 = 4$) with $6$ sewing circles
($3g - 3 = 6$).  Three of these circles correspond to handles
(B-cycles), and the remaining three correspond to internal
seams.

\emph{Step 1: block structure.}
The $3 \times 3$ block operator $\sewker_3$ on
$A^2(D)^{\oplus 3}$ has diagonal entries $\sewop_{q_j}$
for the three handle sewing operators with
$\|\sewop_{q_j}\|_1 = q_j/(1-q_j)$ (trace class), and
off-diagonal entries $R_{ij}$ with
$\|R_{ij}\|_{\hscl} \leq C\,q_{ij}^{1/2}$
(Hilbert--Schmidt) where $q_{ij}$ is the sewing parameter
of the internal seam connecting handles $i$ and $j$.

\emph{Step 2: trace-class verification.}
The operator $\sewker_3$ is trace class on $A^2(D)^{\oplus 3}$.
The trace norm satisfies
\[
\|\sewker_3\|_1
\leq \sum_{j=1}^{3} \|\sewop_{q_j}\|_1
+ \sum_{i \neq j} \|R_{ij}\|_1
\leq \sum_{j=1}^{3} \frac{q_j}{1-q_j}
+ \sum_{i \neq j}
\|R_{ij}\|_{\hscl}^2
< \infty.
\]

\emph{Step 3: Fredholm determinant.}
The Fredholm determinant is computed by the standard $3 \times 3$
block formula:
\begin{align}
\det(1 - \sewker_3)
&= \det(1 - \sewop_{q_1})
\cdot \det(1 - S_{23|1})
\notag\\
&\quad\text{where } S_{23|1}
= \begin{pmatrix} \sewop_{q_2} + R_{21}(1-\sewop_{q_1})^{-1}R_{12}
& R_{23} + R_{21}(1-\sewop_{q_1})^{-1}R_{13}\\
R_{32} + R_{31}(1-\sewop_{q_1})^{-1}R_{12}
& \sewop_{q_3} + R_{31}(1-\sewop_{q_1})^{-1}R_{13}
\end{pmatrix}
\notag
\end{align}
is the Schur complement.  Iterating the Schur complement
reduction gives:
\begin{equation}% label removed: eq:thqg-X-genus-3-iterated-schur
\det(1 - \sewker_3)
= \prod_{j=1}^{3} \det(1 - \sewop_{q_j})
\cdot \det(1 - \Delta_3),
\end{equation}
where $\Delta_3$ is a trace-class operator encoding the
inter-handle coupling.  In the degeneration limit where all
internal seams are pinched ($R_{ij} \to 0$), $\Delta_3 \to 0$
and the determinant factors into a product of three genus-$1$
determinants:
\[
\det(1 - \sewker_3)
\to
\prod_{j=1}^{3}
\dedeta(\tau_j)^2.
\]

\emph{Step 4: Polyakov formula.}
The Polyakov--Alvarez formula at genus~$3$ gives
$\det(1-\sewker_3)^{-1}
= (\det\operatorname{Im}\Omega)^{-1/2}
\cdot (\zetareg\Delta_{\Sigma_3})^{-1/2}$.
\end{proof}

\begin{remark}[The Schottky problem at genus~$3$]%
% label removed: rem:thqg-X-schottky%
\index{Schottky problem}%
At genus~$3$, the Siegel upper half-space $\Siegel_3$ has
complex dimension $\binom{3+1}{2} = 6$, and the moduli space
$\cM_3$ has dimension $3g-3 = 6$.  The Torelli map
$\cM_3 \to \Siegel_3/\Sp_6(\Z)$ is therefore generically
surjective at genus~$3$ (this is the classical fact that the
Schottky locus equals the full Siegel space at genus~$\leq 3$).
At genus~$4$, the Siegel space has dimension $10$ while $\cM_4$
has dimension $9$, and the Schottky locus becomes a proper
subvariety.  The Fredholm determinant formula
\eqref{eq:thqg-X-genus-3-det} is therefore a function on all of
$\Siegel_3$ at genus~$3$, but only on the Schottky locus of
$\Siegel_4$ at genus~$4$.
\end{remark}

\subsubsection{General genus: the recursive structure}

\begin{theorem}[General-genus Fredholm structure;
\ClaimStatusProvedHere]%
% label removed: thm:thqg-X-general-genus-fredholm%
\index{Fredholm determinant!general genus}%
For every genus $g \geq 1$, the partition function of
$\cH_\kappa$ on $\Sigma_g$ with period matrix
$\Omega \in \Siegel_g$ is
\begin{equation}% label removed: eq:thqg-X-general-genus-formula
Z_g(\cH_\kappa;\,\Omega)
\;=\;
\bigl(\det\operatorname{Im}\Omega\bigr)^{-\kappa/2}
\cdot
\bigl(\zetareg\Delta_{\Sigma_g}\bigr)^{-\kappa/2}
\;=\;
\det\bigl(1 - \sewker_g\bigr)^{-\kappa},
\end{equation}
where $\sewker_g$ is a trace-class operator on
$A^2(D)^{\oplus g}$, the genus-$g$ sewing kernel of
Definition~\textup{\ref{def:thqg-X-genus-g-sewing}}.
The following recursive properties hold:
\begin{enumerate}[label=\textup{(\roman*)}]
\item \emph{Separating degeneration.}
  If $\Sigma_g$ degenerates into $\Sigma_{g_1} \cup \Sigma_{g_2}$
  with $g_1 + g_2 = g$, then
  \begin{equation}% label removed: eq:thqg-X-separating-factorization
  Z_g(\cH_\kappa) \to
  Z_{g_1}(\cH_\kappa) \cdot Z_{g_2}(\cH_\kappa).
  \end{equation}
\item \emph{Non-separating degeneration.}
  If $\Sigma_g$ degenerates by pinching a non-separating cycle,
  then
  \begin{equation}% label removed: eq:thqg-X-nonseparating-limit
  Z_g(\cH_\kappa)
  \;\to\;
  \operatorname{Tr}_{A^2}\bigl(\sewop_q \cdot Z_{g-1}(\cH_\kappa)
  \bigr),
  \end{equation}
  where the trace is over the one-particle space of the
  collapsing handle and $q \to 0$ is the sewing parameter.
\item \emph{Free energy.}
  The genus-$g$ free energy
  $F_g(\cH_\kappa) = \log Z_g(\cH_\kappa)$ satisfies
  \begin{equation}% label removed: eq:thqg-X-free-energy-genus
  F_g(\cH_\kappa)
  \;=\;
  -\kappa\,\operatorname{Tr}\log(1 - \sewker_g)
  \;=\;
  \kappa\sum_{k=1}^{\infty}
  \frac{1}{k}\operatorname{Tr}(\sewker_g^k).
  \end{equation}
\end{enumerate}
\end{theorem}

\begin{proof}
The general formula follows from the Heisenberg sewing theorem
(Theorem~\ref{thm:thqg-X-heisenberg-sewing-full}\,(IV)) and
the Polyakov--Alvarez formula at genus~$g$.

Part~(i): in the separating degeneration, the genus-$g$
sewing kernel $\sewker_g$ becomes block-diagonal:
$\sewker_g \to \sewker_{g_1} \oplus \sewker_{g_2}$
as the separating seam parameter $q_s \to 0$.
Hence $\det(1 - \sewker_g) \to
\det(1 - \sewker_{g_1}) \cdot \det(1 - \sewker_{g_2})$.

Part~(ii): in the non-separating degeneration, one diagonal
block $\sewop_{q_j} \to 0$ as $q_j \to 0$, and the Schur
complement formula reduces the rank of the sewing kernel by one.

Part~(iii): $\log\det(1-\sewker_g)^{-\kappa}
= -\kappa\,\operatorname{Tr}\log(1-\sewker_g)
= \kappa\sum_{k} k^{-1}\operatorname{Tr}(\sewker_g^k)$.
\end{proof}

\begin{proposition}[Free energy: $\hat{A}$-genus verification;
\ClaimStatusProvedHere]%
% label removed: prop:thqg-X-free-energy-ahat%
\index{free energy!$\hat{A}$-genus}%
For $\cH_\kappa$ at level $\kappa$, the Faber--Pandharipande
free energy formula
$F_g = \kappa \cdot \lambda_g^{\mathrm{FP}}$ is verified
through $g = 3$:
\begin{center}
\begin{tabular}{c|c|c|c}
$g$ & $\lambda_g^{\mathrm{FP}}$ & $F_g/\kappa$ & Source \\
\hline
$1$ & $\frac{|B_2|}{2 \cdot 2!} = \frac{1}{24}$ & $\frac{1}{24}$
& $\dedeta(\tau)^{-1}$: constant term \\[4pt]
$2$ & $\frac{2^3-1}{2^3}\frac{|B_4|}{4!}
= \frac{7}{8}\cdot\frac{1}{720}
= \frac{7}{5760}$ & $\frac{7}{5760}$
& Polyakov at $g = 2$ \\[4pt]
$3$ & $\frac{2^5-1}{2^5}\frac{|B_6|}{6!}
= \frac{31}{32}\cdot\frac{1}{30240}
= \frac{31}{967680}$ & $\frac{31}{967680}$
& Polyakov at $g = 3$
\end{tabular}
\end{center}
These match the $\hat{A}$-genus coefficients
$\hat{A}_g = \frac{(2^{2g-1}-1)|B_{2g}|}{2^{2g-1}(2g)!}$
(Theorem~\textup{\ref{thm:family-index}}).
\end{proposition}

\begin{proof}
The genus-$g$ free energy of a single free boson is
$F_g = -\frac{1}{2}\log\zetareg\Delta_{\Sigma_g}$
integrated over $\overline{\cM}_g$.  By Faltings'
arithmetic Riemann--Roch formula and the Mumford relations,
$\int_{\overline{\cM}_g} \lambda_g
= \frac{(2^{2g-1}-1)|B_{2g}|}{2^{2g-1}(2g)!}$
(Faber's formula).  The claim follows by direct substitution.
\end{proof}

% =====================================================================
%
%  SECTION 4: SCALAR SATURATION AND CLASS-G ALGEBRAS
%
% =====================================================================

\subsection{Extension to class-G algebras via scalar saturation}
% label removed: subsec:thqg-X-class-G

The Heisenberg sewing theorem gives a Fredholm determinant
formula for the simplest possible chiral algebra, the free boson.
The shadow depth classification
(Definition~\ref{def:shadow-depth-classes}
in Appendix~\ref*{V1-app:nonlinear-modular-shadows})
organizes chiral algebras into four classes based on the
termination behavior of the shadow obstruction tower:
\begin{itemize}
\item Class~\textbf{G} (Gaussian): $r_{\max} = 2$.
  The shadow obstruction tower terminates at weight~$2$:
  $\Theta_\cA^{\min} = \eta\otimes\Gamma_\cA$.
  On the proved uniform-weight scalar lane this specializes to
  $\Theta_\cA^{\min} = \kappa(\cA)\cdot\eta \otimes \Lambda$.
  Examples: $\cH_\kappa$, lattice algebras $V_\Lambda$.
\item Class~\textbf{L} (Lie): $r_{\max} = 3$.
  Examples: $V_k(\fg)$ at generic level.
\item Class~\textbf{C} (contact): $r_{\max} = 4$.
  Examples: $\beta\gamma$ system.
\item Class~\textbf{M} (mixed): $r_{\max} = \infty$.
  Examples: $\mathrm{Vir}_c$, $\cW_N$ for $N \geq 3$.
\end{itemize}
For class-G algebras, the scalar rigidity theorem
(Theorem~\ref{thm:algebraic-family-rigidity}) ensures only that the
minimal MC class is purely scalar in arity.  This does not by itself
identify the full nonperturbative partition function with a Fredholm
determinant.  The Fredholm formula below is proved for Gaussian
algebras on the uniform-weight scalar lane.

\begin{theorem}[Fredholm scalar partition function on the Gaussian
scalar lane; \ClaimStatusProvedHere]%
% label removed: thm:thqg-X-class-G-fredholm%
\index{partition function!class G}%
\index{scalar saturation!Fredholm determinant}%
Let $\cA$ be a class-G Gaussian chiral algebra on the proved
uniform-weight scalar lane
\textup{(}shadow depth $r_{\max} = 2$\textup{)} with
modular characteristic $\kappa = \kappa(\cA)$ and central
charge $c(\cA)$.  Assume HS-sewing
\textup{(}Definition~\textup{\ref{def:hs-sewing}}\textup{)}.
Then the scalar genus-$g$ partition function is a Fredholm
determinant:
\begin{equation}% label removed: eq:thqg-X-class-G-formula
Z_g^{\mathrm{scal}}(\cA;\,\Omega)
\;=\;
Z_g^{\mathrm{G}}(\kappa;\,\Omega)
\;:=\;
\bigl(\det\operatorname{Im}\Omega\bigr)^{-\kappa/2}
\cdot
\bigl(\zetareg\Delta_{\Sigma_g}\bigr)^{-\kappa/2},
\end{equation}
up to a $\kappa$-dependent prefactor encoding the vacuum energy
$q^{-c/24}$.  This formula depends on $\cA$ only through
$\kappa(\cA)$ on that scalar lane.  For general multi-weight
class-G algebras, the identification of the scalar coefficient with
$\kappa(\cA)\Lambda$ remains open, so no full Fredholm theorem is
claimed here.
\end{theorem}

\begin{proof}
On the proved uniform-weight Gaussian lane,
Theorem~\ref{thm:theta-direct-derivation} identifies the minimal
scalar package as
$\Theta_\cA^{\min} = \kappa(\cA)\eta\otimes\Lambda$.
The scalar genus-$g$ partition function is therefore the
$\kappa(\cA)$-multiple of the Heisenberg scalar partition function.
The HS-sewing condition ensures convergence of the Fredholm
determinant (by Theorem~\ref{thm:general-hs-sewing}, which
requires polynomial OPE growth and subexponential sector growth;
class-G algebras satisfy both).

The Fredholm determinant formula then follows from the
Heisenberg computation
(Theorem~\ref{thm:thqg-X-heisenberg-sewing-full}\,(IV))
with $\kappa$ replaced by $\kappa(\cA)$.
\end{proof}

\begin{example}[Lattice vertex algebras]%
% label removed: ex:thqg-X-lattice-fredholm%
\index{lattice algebra!Fredholm determinant}%
Let $V_\Lambda$ be a lattice vertex algebra associated to an
even integral lattice $\Lambda$ of rank $r$.  Then
$\kappa(V_\Lambda) = r$ (the rank of the lattice), and the
genus-$g$ partition function is
\begin{equation}% label removed: eq:thqg-X-lattice-formula
Z_g(V_\Lambda;\,\Omega)
\;=\;
\bigl(\det\operatorname{Im}\Omega\bigr)^{-r/2}
\cdot
\bigl(\zetareg\Delta_{\Sigma_g}\bigr)^{-r/2}
\cdot
\Theta_\Lambda^{(g)}(\Omega),
\end{equation}
where $\Theta_\Lambda^{(g)}(\Omega)
= \sum_{\lambda \in \Lambda^g}
\exp(\pi i \lambda^T \Omega \lambda)$ is the genus-$g$
lattice theta function.  The Fredholm part (first two factors)
carries the oscillator modes and is universal in $r$;
the theta function carries the zero-mode sum over the
lattice and is lattice-specific.

The HS-sewing condition is verified: $V_\Lambda$ has sector
growth $\dim (V_\Lambda)_n \leq |\Lambda^*/\Lambda| \cdot p(n)^r$
(polynomial times partition function growth), and OPE
coefficients are bounded by $\kappa = r$ times the lattice
norm (Corollary~\ref{cor:hs-sewing-standard-landscape}).
\end{example}

\begin{remark}[Curvature-braiding orthogonality for lattice algebras]%
% label removed: rem:thqg-X-lattice-curvature-braiding%
\index{lattice algebra!curvature-braiding orthogonality}%
By Theorem~\ref{thm:lattice:curvature-braiding-orthogonal},
$\kappa(V_\Lambda^{N,q}) = \mathrm{rank}(\Lambda)$
independent of the cocycle~$q$.  This means the Fredholm
determinant part of the partition function is insensitive to
the choice of cocycle; only the theta function
$\Theta_\Lambda^{(g)}$ depends on the lattice embedding.
This orthogonality between the oscillator (Fredholm) sector
and the zero-mode (theta) sector is a direct consequence of
the class-G property.
\end{remark}

\begin{proposition}[Rank additivity;
\ClaimStatusProvedHere]%
% label removed: prop:thqg-X-rank-additivity%
\index{Fredholm determinant!rank additivity}%
If $\cA = \cA_1 \otimes \cA_2$ is a tensor product of
class-G algebras, then
\begin{equation}% label removed: eq:thqg-X-rank-additivity
Z_g(\cA;\,\Omega)
= Z_g(\cA_1;\,\Omega)
\cdot Z_g(\cA_2;\,\Omega),
\end{equation}
reflecting the additivity of the modular characteristic
$\kappa(\cA) = \kappa(\cA_1) + \kappa(\cA_2)$
and the multiplicativity of the Fredholm determinant
$\det(1 - T_1 \oplus T_2)
= \det(1 - T_1) \cdot \det(1 - T_2)$.
\end{proposition}

\begin{proof}
By Proposition~\ref{prop:independent-sum-factorization},
the modular characteristic is additive:
$\kappa(\cA_1 \otimes \cA_2)
= \kappa(\cA_1) + \kappa(\cA_2)$.
The sewing operator for a tensor product is the direct sum:
$\sewker_g(\cA_1 \otimes \cA_2)
= \sewker_g(\cA_1) \oplus \sewker_g(\cA_2)$
on $A^2(D)^{\oplus g_1} \oplus A^2(D)^{\oplus g_2}$.
The Fredholm determinant of a direct sum is the product:
$\det(1 - T_1 \oplus T_2) = \det(1-T_1)\det(1-T_2)$.
\end{proof}

% =====================================================================
%
%  SECTION 5: NON-FREDHOLM CORRECTIONS FOR CLASSES L, C, M
%
% =====================================================================

\subsection{Non-Fredholm corrections for classes L, C, M}
% label removed: subsec:thqg-X-non-fredholm

For algebras of shadow depth $r_{\max} > 2$, the partition
function is no longer a pure Fredholm determinant.  The higher
shadow data ($\fC$, $\mathfrak{Q}$, and all higher
$o_r$ for $r \geq 3$) produce perturbative corrections
to the Gaussian (Fredholm) base, organized as Feynman
integrals weighted by the graph-sum formula of
Definition~\ref*{V1-def:modular-bar-hamiltonian}.

\begin{definition}[Non-Fredholm correction]%
% label removed: def:thqg-X-non-fredholm-correction%
\index{non-Fredholm correction|textbf}%
\index{partition function!non-Fredholm correction}%
For a chiral algebra $\cA$ with shadow depth $r_{\max}$, define
the \emph{genus-$g$ non-Fredholm correction} as
\begin{equation}% label removed: eq:thqg-X-non-fredholm-correction
\delta Z_g(\cA;\,\Omega)
\;:=\;
Z_g(\cA;\,\Omega)
\;-\;
Z_g^{\mathrm{G}}(\kappa(\cA);\,\Omega),
\end{equation}
where $Z_g^{\mathrm{G}}$ is the Gaussian (Fredholm) base of
\eqref{eq:thqg-X-class-G-formula}.  The correction
$\delta Z_g$ encodes the contribution of all shadow data
beyond the modular characteristic $\kappa$.
\end{definition}

\begin{theorem}[Feynman integral expansion;
\ClaimStatusProvedHere]%
% label removed: thm:thqg-X-feynman-expansion%
\index{Feynman integral!non-Fredholm correction}%
For a modular Koszul chiral algebra $\cA$ satisfying HS-sewing,
the genus-$g$ non-Fredholm correction admits a Feynman integral
expansion:
\begin{equation}% label removed: eq:thqg-X-feynman-expansion
\delta Z_g(\cA;\,\Omega)
\;=\;
\sum_{r=3}^{r_{\max}}
\sum_{\Gamma \in \mathcal{G}^{\mathrm{conn}}_{g,r}}
\frac{1}{|\Aut(\Gamma)|}
\int_{\Sigma_g^{v(\Gamma)}}
\omega_\Gamma(\Omega),
\end{equation}
where:
\begin{itemize}
\item $\mathcal{G}^{\mathrm{conn}}_{g,r}$ is the set of
  connected stable graphs of genus~$g$ with vertices of
  arity at most~$r$;
\item $v(\Gamma)$ is the number of vertices of~$\Gamma$;
\item $\omega_\Gamma(\Omega)$ is the integrand determined by
  the graph~$\Gamma$, with vertex labels from the transferred
  cyclic minimal model $\ell_k^{\mathrm{tr}}$ of $\cA$ and
  edge contractions by the complementarity propagator
  $P_\cA = H_\cA^{-1}$ (the inverse of the Hodge Laplacian).
\end{itemize}
The Feynman expansion converges absolutely for $|q_\alpha| < 1$
(all sewing parameters inside the unit disk).
\end{theorem}

\begin{proof}
The genus-$g$ partition function is computed by the shadow
Postnikov tower: $Z_g = \langle e^{\Theta_\cA^{(g)}} \rangle_g$.
The Gaussian part $Z_g^{\mathrm{G}}$ is the contribution of
$\Theta_\cA^{\leq 2}$, which consists of the modular
characteristic $\kappa$ contracted with the Hodge cycle.
The non-Fredholm correction is the contribution of
$\Theta_\cA - \Theta_\cA^{\leq 2}$, which involves the
higher-arity terms $\ell_k^{\mathrm{tr}}$ for $k \geq 3$.

Expanding $e^{\Theta}$ in the graph-sum formula
(Definition~\ref*{V1-def:modular-bar-hamiltonian}), each
graph $\Gamma$ with a vertex of arity $r > 2$ contributes to
$\delta Z_g$.  The weight of a graph is
$w(\Gamma) = \sum_{v \in V(\Gamma)} \mathrm{arity}(v)$,
and graphs with $w > 2|V(\Gamma)|$ contribute to the
non-Fredholm correction.

Convergence: the HS-sewing condition
(Theorem~\ref{thm:general-hs-sewing}) ensures that each
Feynman integral converges absolutely.  The propagator
$P_\cA = H_\cA^{-1}$ on $\Sigma_g$ has the same exponential
decay as the Bergman composition operator
(Lemma~\ref{lem:thqg-X-composition-decay}), and the vertex
labels $\ell_k^{\mathrm{tr}}$ are bounded by the polynomial
OPE growth.  The arity-cutoff lemma (Lemma~\ref{lem:arity-cutoff})
bounds the number of contributing graphs at each order.
\end{proof}

\subsubsection{Class L: cubic corrections from Lie algebras}

\begin{proposition}[Class-L Feynman integrals;
\ClaimStatusProvedHere]%
% label removed: prop:thqg-X-class-L-feynman%
\index{class L!Feynman integrals}%
For a class-L algebra $\cA$ \textup{(}$r_{\max} = 3$,
e.g.\ $V_k(\fg)$ at generic level\textup{)}, the
non-Fredholm correction at genus~$g$ involves only
trivalent graphs:
\begin{equation}% label removed: eq:thqg-X-class-L-correction
\delta Z_g^{\mathrm{L}}(\cA;\,\Omega)
\;=\;
\sum_{\Gamma \in \mathcal{G}^{3\text{-val}}_{g}}
\frac{1}{|\Aut(\Gamma)|}
\int_{\Sigma_g^{v(\Gamma)}}
\omega_\Gamma^{(3)}(\Omega),
\end{equation}
where $\mathcal{G}^{3\text{-val}}_g$ denotes the set of
connected $3$-valent stable graphs of genus~$g$, and
$\omega_\Gamma^{(3)}$ is the integrand with cubic vertex
labels $\ell_3^{\mathrm{tr}}$ (the transferred $L_\infty$
ternary bracket).

A $3$-valent graph of genus~$g$ with $v$ vertices has
$e = 3v/2$ edges and loop number $g = e - v + 1 = v/2 + 1$,
so $v = 2(g-1)$ and $e = 3(g-1)$.  In particular:
\begin{center}
\begin{tabular}{c|c|c|c}
$g$ & $v$ & $e$ & \# connected $3$-valent graphs \\
\hline
$1$ & $0$ & $0$ & $0$ (no correction) \\
$2$ & $2$ & $3$ & $2$ (theta, sunset) \\
$3$ & $4$ & $6$ & \text{finitely many}
\end{tabular}
\end{center}
The first non-trivial correction appears at genus~$2$
with the theta graph and the sunset graph.
\end{proposition}

\begin{proof}
For class-L algebras, $\ell_k^{\mathrm{tr}} = 0$ for $k \geq 4$
(the shadow obstruction tower terminates at arity~$3$).  Hence only graphs
with vertices of arity exactly~$3$ contribute to $\delta Z_g$.
These are the $3$-regular (trivalent) graphs.

The graph count follows from the Euler relation: a connected
graph with $v$ vertices, $e$ edges, and first Betti number
$g$ satisfies $g = e - v + 1$.  For $3$-regular graphs,
$3v = 2e$ (every edge has two endpoints; every vertex has
three), so $e = 3v/2$ and $g = v/2 + 1$.
\end{proof}

\begin{example}[Affine $\AffKM{g}_k$ at genus~$2$]%
% label removed: ex:thqg-X-affine-genus-2%
\index{affine Kac--Moody!genus-2 correction}%
For $V_k(\fg)$ at generic level $k \neq -h^\vee$, the
genus-$2$ non-Fredholm correction involves two Feynman graphs.
Let $f_{abc}$ denote the structure constants of $\fg$
and $P(z,w;\Omega)$ the genus-$2$ propagator.  Then:
\begin{enumerate}[label=\textup{(\alph*)}]
\item \emph{Theta graph}
  $(\Gamma_\theta$: two vertices, three edges$)$:
  \begin{equation}% label removed: eq:thqg-X-theta-graph
  I_{\theta}
  = \frac{1}{12}
  f_{abc}\,f_{a'b'c'}\,
  \langle a,a'\rangle\langle b,b'\rangle\langle c,c'\rangle
  \int_{\Sigma_2 \times \Sigma_2}
  P(z,w;\Omega)^3\,dz\,dw.
  \end{equation}
\item \emph{Sunset graph}
  $(\Gamma_{\mathrm{sun}}$: two vertices, three edges,
  one edge connects vertex to itself$)$:
  \begin{equation}% label removed: eq:thqg-X-sunset-graph
  I_{\mathrm{sun}}
  = \frac{1}{4}
  f_{abc}\,f_{a'b'c'}\,
  \langle a,a'\rangle
  \int_{\Sigma_2}
  \langle b, c\rangle_z
  \,P(z,w;\Omega)\,
  \langle b', c'\rangle_w\,
  dz\,dw.
  \end{equation}
\end{enumerate}
The total correction is
$\delta Z_2^{\mathrm{L}}(V_k(\fg);\,\Omega)
= I_\theta + I_{\mathrm{sun}}$.
At integrable level $k \in \Z_{\geq 0}$, these Feynman
integrals reduce to combinatorial sums via the Verlinde
formula (Example~\ref{eq:verlinde-general}).

% ----------------------------------------------------------------
% GENUS-2 FEYNMAN GRAPHS FOR CLASS L
% ----------------------------------------------------------------
\begin{center}
\begin{tikzpicture}[x=1cm, y=1cm]
% ---- Theta graph ----
\node[font=\footnotesize\scshape] at (0,1.2) {$\Gamma_\theta$};
\node[gv0, label={[font=\tiny]below left:$f_{abc}$}]
  (t1) at (-0.4,0) {};
\node[gv0, label={[font=\tiny]below right:$f_{a'b'c'}$}]
  (t2) at (0.4,0) {};
\draw[iedge] (t1) to[bend left=40] node[font=\tiny, above]
  {$P$} (t2);
\draw[iedge] (t1) -- (t2);
\draw[iedge] (t1) to[bend right=40] (t2);
\node[font=\tiny, text=black!50] at (0,-0.8) {$I_\theta
  = \frac{1}{12}\!\int P(z,w)^3$};
%
% ---- Plus sign ----
\node[font=\normalsize] at (2.2,0) {$+$};
%
% ---- Sunset graph ----
\node[font=\footnotesize\scshape] at (4.2,1.2) {$\Gamma_{\mathrm{sun}}$};
\node[gv0, label={[font=\tiny]below left:$f_{abc}$}]
  (s1) at (3.8,0) {};
\node[gv0, label={[font=\tiny]below right:$f_{a'b'c'}$}]
  (s2) at (4.6,0) {};
\draw[iedge] (s1) -- (s2);
\draw[iedge] (s1) to[bend left=40] (s2);
\draw[sloop] (s1) to[out=150,in=210,looseness=5] (s1);
\node[font=\tiny, text=black!50] at (4.2,-0.8)
  {$I_{\mathrm{sun}} = \frac{1}{4}\!\int
    \langle b,c\rangle_z\, P\, \langle b',c'\rangle_w$};
%
% ---- Equals ----
\node[font=\normalsize] at (6.6,0) {$=$};
\node[font=\footnotesize] at (8.2,0)
  {$\delta Z_2^{\mathrm{L}}(V_k(\fg);\,\Omega)$};
\end{tikzpicture}
\end{center}
\end{example}

\subsubsection{Class C: quartic contact corrections}

\begin{proposition}[Class-C quartic contact;
\ClaimStatusProvedHere]%
% label removed: prop:thqg-X-class-C-quartic%
\index{class C!quartic contact}%
For a class-C algebra $\cA$ with $r_{\max} = 4$
\textup{(}e.g.\ the $\beta\gamma$ system\textup{)},
the non-Fredholm correction at genus~$g$ involves
graphs with vertices of arity $3$ and $4$.
The quartic vertex $\ell_4^{\mathrm{tr}}$ is governed by the
quartic resonance class
$\mathfrak{Q} \in H^2(\cA^{\mathrm{sh}}_{4,0})$
of Appendix~\textup{\ref*{V1-app:nonlinear-modular-shadows}}.

The first quartic contribution appears at genus~$1$:
\begin{equation}% label removed: eq:thqg-X-quartic-genus-1
\delta Z_1^{(4)}(\cA;\,\tau)
\;=\;
\mathfrak{Q}(\cA)\,
\int_{E_\tau}
P(z;\tau)^2\,dz,
\end{equation}
where $P(z;\tau) = G_\tau(z)$ is the genus-$1$ propagator.
For the $\beta\gamma$ system,
$\mathfrak{Q}_{\beta\gamma} = 0$
\textup{(}Corollary~\textup{\ref*{V1-cor:nms-betagamma-mu-vanishing}}\textup{)},
so the quartic contact invariant vanishes and the first
genuine quartic correction appears at genus~$2$.
\end{proposition}

\begin{proof}
The quartic vertex contributes Feynman graphs with at least one
$4$-valent vertex.  At genus~$1$, the simplest such graph is a
single $4$-valent vertex with two self-loops, contributing
$\delta Z_1^{(4)}$.  The quartic resonance class $\mathfrak{Q}$
is the weight-$4$ component of the shadow algebra $\cA^{\mathrm{sh}}$
(Corollary~\ref*{V1-cor:shadow-extraction}).

For the $\beta\gamma$ system: the quartic contact invariant
$\mu_{\beta\gamma} = \langle \eta, m_3(\eta,\eta,\eta)\rangle = 0$
by Corollary~\ref*{V1-cor:nms-betagamma-mu-vanishing} (rank-one abelian
rigidity on the weight-changing line).
\end{proof}

\subsubsection{Class M: the infinite shadow obstruction tower}

\begin{remark}[Class-M non-Fredholm corrections]%
% label removed: rem:thqg-X-class-M-corrections%
\index{class M!non-Fredholm corrections}%
For class-M algebras ($r_{\max} = \infty$, e.g.\ $\mathrm{Vir}_c$
and $\cW_N$ for $N \geq 3$), the Feynman expansion involves
vertices of arbitrarily high arity.  The genus-$g$ non-Fredholm
correction is an infinite series:
\[
\delta Z_g^{\mathrm{M}}(\cA;\,\Omega)
= \sum_{r=3}^{\infty}
\sum_{\Gamma \in \mathcal{G}^{\mathrm{conn}}_{g,r}}
\frac{1}{|\Aut(\Gamma)|}
\int_{\Sigma_g^{v(\Gamma)}}
\omega_\Gamma(\Omega).
\]
Convergence of this series is guaranteed by the HS-sewing
condition (Theorem~\ref{thm:general-hs-sewing}), which provides
uniform bounds on the vertex labels $\ell_k^{\mathrm{tr}}$, and
the arity-cutoff lemma (Lemma~\ref{lem:arity-cutoff}), which
ensures that the contribution of high-arity graphs is
exponentially suppressed.

For the Virasoro algebra at genus~$1$, the non-Fredholm
correction involves the quartic contact invariant
$\mathfrak{Q}_{\mathrm{Vir}} = 10/[c(5c+22)]$
and all higher obstruction classes $o_r$ for $r \geq 5$
(the quintic obstruction $o_5$ is forced to be nonzero
by Theorem~\ref{thm:quintic-forced}).  The full partition
function $Z_1(\mathrm{Vir}_c;\,\tau)$ at genus~$1$ is:
\[
Z_1(\mathrm{Vir}_c;\,\tau)
= \dedeta(\tau)^{-(c-1)}
= q^{-(c-1)/24}\prod_{n=2}^{\infty}(1-q^n)^{-1},
\]
the character of the Virasoro vacuum module, which is
\emph{not} a pure Fredholm determinant on a single Bergman space
(it is missing the $n=1$ factor from the product).
\end{remark}

\begin{proposition}[Virasoro: Fredholm base + corrections;
\ClaimStatusProvedHere]%
% label removed: prop:thqg-X-virasoro-decomposition%
\index{Virasoro!Fredholm decomposition}%
The genus-$1$ Virasoro partition function admits the decomposition
\begin{equation}% label removed: eq:thqg-X-virasoro-decomposition
Z_1(\mathrm{Vir}_c;\,\tau)
\;=\;
\underbrace{|\dedeta(\tau)|^{-2c}}_{\text{Gaussian base}
(\kappa = c/2)}
\;\cdot\;
\underbrace{|\dedeta(\tau)|^{2}
\cdot (1-q)}_{\text{null-state correction}},
\end{equation}
where the Gaussian base is the Fredholm determinant at
$\kappa = c/2$ and the correction factor removes the
contribution of the $L_{-1}$ null state.
\end{proposition}

\begin{proof}
The Virasoro character at genus~$1$ is
$\chi_c(\tau) = q^{(1-c)/24}/(1-q)\prod_{n=2}^{\infty}(1-q^n)^{-1}
= q^{-c/24}\,\dedeta(\tau)^{-1}\cdot(1-q)^{-1}\cdot\dedeta(\tau)/q^{-1/24}$\,.
More directly:
\begin{align}
\chi_c(\tau)
&= q^{(1-c)/24} \prod_{n=2}^{\infty}(1-q^n)^{-1}
\notag\\
&= q^{(1-c)/24}
\cdot \frac{\prod_{n=1}^{\infty}(1-q^n)^{-1}}{(1-q)^{-1}}
\notag\\
&= q^{(1-c)/24}
\cdot (1-q) \cdot \prod_{n=1}^{\infty}(1-q^n)^{-1}
\notag\\
&= (1-q)\cdot q^{-c/24}\cdot q^{1/24}
\cdot \prod_{n=1}^{\infty}(1-q^n)^{-1}
\notag\\
&= (1-q)\cdot q^{-c/24}\cdot \dedeta(\tau)^{-1}.
% label removed: eq:thqg-X-virasoro-factorization
\end{align}
The factor $(1-q)$ removes the $L_{-1}$ descendant (which is
null for the vacuum module), and $\dedeta(\tau)^{-1}$ is the
oscillator contribution $\prod_{n=1}^{\infty}(1-q^n)^{-1}$
up to the $q^{1/24}$ prefactor.  The Gaussian base
$|\dedeta|^{-2c}$ arises from $c$ free-field copies at
$\kappa = 1$ each, giving $\kappa_{\mathrm{total}} = c/2$.
\end{proof}

\begin{remark}[Resurgent structure for class-M algebras]%
% label removed: rem:thqg-X-resurgent%
\index{resurgence!class M}%
For class-M algebras, the Feynman expansion
$\delta Z_g^{\mathrm{M}}$ is expected to be resurgent:
the asymptotic (divergent) perturbative series in the shadow
tower is Borel summable, with non-perturbative corrections
controlled by the resonance rank $\rho(\cA)$
(Definition~\ref{def:resonance-rank}).  The Stokes phenomenon
as $\tau$ crosses walls of marginal stability corresponds to
the wall-crossing of the MC moduli of
$\Theta_\cA \in \MC(\gAmod)$.  This resurgent structure is
conjectural beyond the Virasoro case, where the Zamolodchikov
recursion provides an independent control.
\end{remark}

% =====================================================================
%
%  SECTION 6: ANALYTIC ASPECTS
%
% =====================================================================

\subsection{Analytic aspects: the analytic bar coalgebra,
coderived categories, and open problems}
% label removed: subsec:thqg-X-analytic-aspects

\subsubsection{The analytic bar coalgebra}

\begin{definition}[Analytic bar coalgebra; review]%
% label removed: def:thqg-X-analytic-bar-review%
\index{bar coalgebra!analytic!review}%
Recall from Definition~\ref{def:analytic-bar-coalgebra} that the
analytic bar coalgebra $\anbar(\cA)$ is defined by closing
the algebraic bar construction inside the sewing envelope:
\[
\anbar_n(\cA)
\;=\;
\overline{
\algcore^{\otimes n}
\otimes \Omega^*_{\log}(C_n(X))
}^{\,\graphnorm(d_{\mathrm{bar}})},
\]
where the graph norm of $d_{\mathrm{bar}}$ is the norm
$\|x\|_{\mathrm{graph}}^2 := \|x\|^2 + \|d_{\mathrm{bar}} x\|^2$.
This is the analytic closure of the algebraic bar construction
(Definition~\ref*{V1-def:bar-differential-complete}).
\end{definition}

\begin{proposition}[Analytic bar differential is bounded;
\ClaimStatusProvedHere]%
% label removed: prop:thqg-X-analytic-bar-bounded%
\index{bar coalgebra!analytic!boundedness}%
Under HS-sewing \textup{(}Definition~\textup{\ref{def:hs-sewing}}\textup{)},
the bar differential $d_{\mathrm{bar}} \colon \anbar_n(\cA) \to
\anbar_{n-1}(\cA)$ is a bounded operator for the graph-norm
topology, with operator norm
\begin{equation}% label removed: eq:thqg-X-bar-bounded
\|d_{\mathrm{bar}}\|_{\mathrm{op}}
\;\leq\;
\binom{n}{2}\,K\,\sup_q \|m_q\|_{\hscl},
\end{equation}
where $K$ is the polynomial OPE growth constant and $m_q$ is the
$q$-weighted multiplication.
\end{proposition}

\begin{proof}
The bar differential $d_{\mathrm{bar}} = \sum_{i<j}
\operatorname{Res}_{z_i=z_j} \circ Y_{ij}$ has
$\binom{n}{2}$ summands, each of which is a composition of
the OPE operator $Y_{ij}$ (bounded by $K$ on the sewing envelope)
and the residue operator (norm $1$ on the graph-norm closure).
The HS-sewing condition ensures $\|m_q\|_{\hscl} < \infty$,
which controls the operator norm of $Y_{ij}$ on the weighted
completion $\sewenv_q$.
\end{proof}

\begin{proposition}[Analytic coproduct;
\ClaimStatusProvedHere]%
% label removed: prop:thqg-X-analytic-coproduct%
\index{bar coalgebra!analytic!coproduct}%
The analytic bar coalgebra $\anbar(\cA)$ carries a continuous
coproduct
\begin{equation}% label removed: eq:thqg-X-analytic-coproduct
\bar{\Delta}^{\mathrm{an}} \colon
\anbar(\cA)
\;\to\;
\anbar(\cA)
\;\widehat{\otimes}\;
\anbar(\cA),
\end{equation}
where $\widehat{\otimes}$ is the completed tensor product
in the graph-norm topology.  The coproduct extends the algebraic
coproduct $\bar{\Delta}$ of the algebraic bar construction
by continuity.
\end{proposition}

\begin{proof}
The algebraic coproduct $\bar{\Delta}$ on $\barB_n(\cA)$
sends $\barB_n \to \bigoplus_{p+q=n} \barB_p \otimes \barB_q$
by deshuffle.  Each deshuffle component is a restriction map
(restricting meromorphic forms from $C_n$ to
$C_p \times C_q$), which is continuous for the graph-norm
topology.  The completed coproduct is the closure of the
algebraic coproduct.
\end{proof}

\begin{remark}[The bar Laplacian and Hodge decomposition]%
\label{rem:bar-laplacian-hodge}%
\index{bar Laplacian|textbf}%
\index{Hodge decomposition!bar complex}
For a unitary VOA~$\cA$ satisfying~(AU1), the formal adjoint
$d_B^*$ is densely defined on~$B^{\mathrm{an}}(\cA)$
(it is the BPZ-conjugate creation operator, reversing the
collision contraction).  The \emph{bar Laplacian}
$\Delta_B := d_B\,d_B^* + d_B^*\,d_B$
is essentially self-adjoint on the algebraic bar complex, because
each weight space $B_n(\cA)_{\mathrm{weight}=N}$ is
finite-dimensional (a basic VOA axiom), and a symmetric operator
on a finite-dimensional space is automatically self-adjoint.
The full operator
$\Delta_B = \bigoplus_{N,n} \Delta_B|_{N,n}$
is an orthogonal direct sum of finite-dimensional self-adjoint
operators, hence essentially self-adjoint on the algebraic core.

Consequences: (i)~discrete spectrum with finite multiplicities,
accumulating only at~$+\infty$;
(ii)~the bar Hodge theorem
$\ker(\Delta_B) = H^*(B(\cA))$
(harmonic forms represent bar cohomology);
(iii)~for chirally Koszul~$\cA$, bar cohomology concentrates on
the diagonal (weight $=$ bar degree), giving a positive spectral
gap off the diagonal.  Compare the screening Hodge theorem
(Theorem~\ref{thm:screening-hodge-theorem} in
\S\ref{sec:ym-spectral-gap}) for the analogous statement in the
Yang--Mills setting.
\end{remark}

\subsubsection{Coderived categories and genus $g \geq 1$}

\begin{remark}[Curvature forces coderived passage; review]%
% label removed: rem:thqg-X-coderived-review%
\index{coderived category!Fredholm context}%
At genus~$g \geq 1$, the fiberwise bar differential satisfies
$\dfib^{\,2} = \kappa(\cA) \cdot \omega_g \neq 0$
(the curvature of the Hodge connection;
Theorem~\ref{thm:genus-graded-bar}).  This means that the
bar complex at genus~$g$ is \emph{curved}: $d^2 \neq 0$,
and ordinary cohomology is not defined.  The correct
receptacle is the \emph{coderived category}
$\codercat(\anbar_g(\cA)\text{-comod})$
(Definition~\ref{def:analytic-shadow} and
Remark~\ref{rem:curvature-coderived-analytic}), not the
ordinary derived category.

For the Fredholm partition functions of this section, the
coderived passage manifests as follows: the genus-$g$
partition function $Z_g$ is the \emph{supertrace} of the
sewing operator on the coderived category of the analytic bar
coalgebra, not an ordinary cohomological invariant.  The
Fredholm determinant formula
$Z_g = \det(1 - \sewker_g)^{-\kappa}$ computes this
supertrace directly via the one-particle reduction, bypassing
the coderived formalism for the Heisenberg and class-G cases.
For classes L, C, M, the coderived category is essential: the
non-Fredholm corrections $\delta Z_g$ are coderived
invariants that vanish in the ordinary derived category.
\end{remark}

\begin{definition}[Coderived shadow partition function]%
% label removed: def:thqg-X-coderived-shadow-pf%
\index{partition function!coderived shadow}%
For a chiral algebra $\cA$ satisfying HS-sewing, define the
\emph{coderived shadow partition function} at genus~$g$ as
\begin{equation}% label removed: eq:thqg-X-coderived-shadow-pf
Z_g^{\mathrm{co}}(\cA;\,\tau)
\;=\;
\operatorname{Str}_{Q_g^{\mathrm{an}}(\cA)}
\!\bigl(\sigma_g\,q^{L_0 - c/24}\bigr),
\end{equation}
where $Q_g^{\mathrm{an}}(\cA)
= \codercat(\anbar_g(\cA)\text{-comod})$ is the
genus-$g$ analytic shadow
\textup{(}Definition~\textup{\ref{def:analytic-shadow}}\textup{)},
$\sigma_g$ is the complementarity involution, and
$\operatorname{Str}$ is the supertrace.
\end{definition}

\begin{proposition}[Coderived = Fredholm for class~G;
\ClaimStatusProvedHere]%
% label removed: prop:thqg-X-coderived-fredholm-G%
\index{coderived category!Fredholm for class G}%
For class-G algebras, the coderived shadow partition function
coincides with the Fredholm partition function:
\begin{equation}% label removed: eq:thqg-X-coderived-fredholm-agreement
Z_g^{\mathrm{co}}(\cA;\,\tau)
\;=\;
Z_g(\cA;\,\Omega)
\;=\;
\det(1 - \sewker_g)^{-\kappa(\cA)}.
\end{equation}
\end{proposition}

\begin{proof}
For class-G algebras, the curvature is scalar
($\dfib^2 = \kappa \cdot \omega_g$), and the coderived category
reduces to the ordinary derived category twisted by the Hodge
line bundle $\hodgebdle^{\kappa/2}$.  The supertrace on
the twisted derived category equals the Fredholm determinant
of the sewing kernel, by the Heisenberg sewing theorem
(Theorem~\ref{thm:thqg-X-heisenberg-sewing-full}).
\end{proof}

\subsubsection{The analytic Koszul pair}

\begin{remark}[Analytic Koszul duality]%
% label removed: rem:thqg-X-analytic-koszul%
\index{Koszul pair!analytic}%
An analytic Koszul pair $(\cA, \cA^!)$ consists of an
algebraic Koszul pair equipped with compatible sewing
envelopes and analytic bar coalgebras
(Definition~\ref{def:analytic-koszul-pair}).  For the
Heisenberg algebra, the analytic Koszul pair is
$(\cH_\kappa, \cH_\kappa^!)$ where
$\cH_\kappa^! = \Sym^{\mathrm{ch}}(V^*)$ with
$\kappa(\cH_\kappa^!) = -\kappa$
(note: $\cH_\kappa^! \neq \cH_{-\kappa}$; the Koszul dual
is a commutative chiral algebra, not a Heisenberg algebra).
The analytic bar coalgebra of $\cH_\kappa^!$
is $\anbar(\cH_\kappa^!) \cong
\overline{\operatorname{Sym} A^2(D)}^{\mathsf{co}}$, the
cofree coalgebra completion of the Bergman Fock space.

The genus-$g$ complementarity decomposition
$Z_g(\cH_\kappa) + Z_g(\cH_\kappa^!)
= (\det\operatorname{Im}\Omega)^{-\kappa/2}
\cdot (\zetareg\Delta)^{-\kappa/2}
+ (\det\operatorname{Im}\Omega)^{\kappa/2}
\cdot (\zetareg\Delta)^{\kappa/2}$
is a direct Fredholm manifestation of Theorem~C
(complementarity): what $\cH_\kappa$ sees as the positive
curvature $+\kappa \cdot \omega_g$, its Koszul dual $\cH_\kappa^!$
sees as negative curvature $-\kappa \cdot \omega_g$.
\end{remark}

\subsubsection{The platonic chain and the sewing hierarchy}

\begin{remark}[The platonic chain for Fredholm partition functions]%
% label removed: rem:thqg-X-platonic-chain%
\index{platonic ideal form!Fredholm partition functions}%
The platonic chain of Remark~\ref{rem:platonic-ideal-form}
has the following concrete incarnation for Fredholm partition
functions:
\[
\underbrace{\algcore}_{\text{algebraic OPE data}}
\;\subset\;
\underbrace{\sewenv}_{\text{sewing envelope}}
\;\rightsquigarrow\;
\underbrace{\operatorname{Sym} A^2(D)}_{\text{Bergman Fock space}}
\;\rightsquigarrow\;
\underbrace{Z_g = \det(1 - \sewker_g)^{-\kappa}}_{\text{Fredholm
partition function}}.
\]
The first arrow (dense embedding) is the mode--Bergman
correspondence; the second arrow (Hilbert completion) uses
the sewing envelope topology; the third arrow (Fredholm
determinant) uses the one-particle reduction and the
Polyakov formula.
Each stage is necessary:
\begin{itemize}
\item $\algcore$ alone gives only formal power series in~$q$
  (the characters, which converge for $|q| < 1$ but have no
  intrinsic analytic continuation).
\item $\sewenv$ promotes the formal power series to operators
  on a Hilbert space, with explicit spectral data.
\item The Fredholm determinant provides the non-perturbative
  completion: it defines $Z_g$ as a function on $\Siegel_g$
  (the Siegel upper half-space), not merely as a formal
  $q$-series.
\end{itemize}
\end{remark}

\subsubsection{Open problems and the analytic frontier}

\begin{openproblem}[Lattice sewing envelope]%
% label removed: prob:thqg-X-lattice-sewing%
\index{lattice algebra!sewing envelope}%
Construct the sewing envelope of a lattice vertex algebra
$V_\Lambda$ explicitly as a completion of the Bergman Fock
space tensored with the lattice group algebra.
Conjecture~\textup{\ref{conj:lattice-sewing}} predicts that
\begin{equation}% label removed: eq:thqg-X-lattice-sewing-conj
V_\Lambda^{\mathrm{sew}}
\;\cong\;
\operatorname{Sym} A^2(D)^{\otimes r}
\;\otimes\;
\ell^2(\Lambda^*/\Lambda),
\end{equation}
where $r = \mathrm{rank}(\Lambda)$ and
$\ell^2(\Lambda^*/\Lambda)$ is the Hilbert space of the
charge lattice.
\end{openproblem}

%% ---- The lattice Steinberg ----

\subsubsection{The lattice Steinberg}
% label removed: sec:lattice-steinberg
\index{Steinberg object!lattice}%
\index{lattice algebra!Steinberg object}%

\providecommand{\Steinberg}{\mathfrak{S}}

\begin{construction}[Lattice Steinberg object]%
% label removed: con:lattice-steinberg%
For a lattice vertex algebra $V_\Lambda$ with $\Lambda$ an
even lattice of rank~$r$, define the \emph{lattice Steinberg
object}
\begin{equation}% label removed: eq:lattice-steinberg-object
\Steinberg_\Lambda
\;=\;
\bigl\{(x_i, \lambda_i)
  \in \FM_k(\bC) \times \Lambda^k
  : \textstyle\sum_i \lambda_i = 0\bigr\}
\;\times\;
\Conf_k^{<}(\bR).
\end{equation}
The constraint $\sum_i \lambda_i = 0$ is charge conservation:
the total lattice momentum flowing through a bar element of
degree~$k$ vanishes.  The Borel--Moore homology decomposes by
coset of the discriminant group:
\begin{equation}% label removed: eq:lattice-steinberg-BM
H_*^{\mathrm{BM}}(\Steinberg_\Lambda)
\;=\;
\bigoplus_{[\mu] \in \Lambda^*/\Lambda}
H_*^{\mathrm{BM}}\!\bigl(\Steinberg_\Lambda^{[\mu]}\bigr),
\end{equation}
where $\Steinberg_\Lambda^{[\mu]}$ restricts the lattice
labels to the coset~$[\mu]$.
\end{construction}

\begin{conjecture}[Conditional on
Conjecture~\textup{\ref{conj:lattice-sewing}};
\ClaimStatusConjectured]%
% label removed: prop:lattice-steinberg-sewing%
Assume the lattice sewing envelope conjecture
\textup{(}Conjecture~\textup{\ref{conj:lattice-sewing})}.
Then the sewing envelope
$\operatorname{Sym} A^2(D)^{\otimes r}
\otimes \ell^2(\Lambda^*/\Lambda)$
is the $L^2$-completion of
$H_*^{\mathrm{BM}}(\Steinberg_\Lambda)$.
The factor $\ell^2(\Lambda^*/\Lambda)$ is the space of
functions on the component group
$\pi_0(\Steinberg_\Lambda) = \Lambda^*/\Lambda$
(the discriminant form).
\end{conjecture}

\begin{proof}
The FM$\times$Conf product decomposition of
$\Steinberg_\Lambda$ gives a direct-sum splitting indexed by
$\Lambda^*/\Lambda$.  On each component the continuous part is
$\FM_k(\bC) \times \Conf_k^{<}(\bR)$ with the charge-zero
sublattice fibre; its BM homology is the $k$-th bar piece of
$\operatorname{Sym} A^2(D)^{\otimes r}$ by the Bergman--Fock
identification of~\S\ref{sec:bergman-fock}.  Summing over
components produces the tensor factor
$\ell^2(\Lambda^*/\Lambda)$, and the $L^2$-completion follows
from the sewing-envelope topology
of~\eqref{eq:thqg-X-lattice-sewing-conj}.
\end{proof}

\begin{remark}[Non-unimodular extension]%
% label removed: rem:lattice-steinberg-non-unimodular%
When $\Lambda$ is not unimodular, $\Lambda^*/\Lambda$ is
nontrivial and contributes additional Stokes data.  The
discriminant form carries a quadratic refinement
$q \colon \Lambda^*/\Lambda \to \bQ/\bZ$ that controls the
phases in the modular transformation; this is the additional
data needed for the non-unimodular sewing.
\end{remark}

\begin{openproblem}[Boundary bar duality]%
% label removed: prob:thqg-X-boundary-bar%
\index{bar coalgebra!boundary duality}%
Establish a Fredholm-level analytic Koszul duality:
for every analytic Koszul pair $(\cA, \cA^!)$ satisfying
HS-sewing, construct a Fredholm equivalence
\begin{equation}% label removed: eq:thqg-X-boundary-bar-conj
\codercat(\anbar(\cA)\text{-comod})
\;\simeq\;
\contrcat(\cA^!\text{-mod})
\end{equation}
between the coderived category of analytic bar comodules and
the contraderived category of $\cA^!$-modules.
This would be the analytic avatar of bar-cobar inversion
(Theorem~B) at genus~$g \geq 1$.
\end{openproblem}

\begin{openproblem}[Non-Fredholm corrections: closed-form expressions]%
% label removed: prob:thqg-X-non-fredholm-closed%
\index{non-Fredholm correction!closed-form}%
For class-L algebras (e.g.\ $V_k(\fg)$), express the
genus-$g$ non-Fredholm correction $\delta Z_g^{\mathrm{L}}$
in terms of Siegel modular forms and conformal blocks.
At integrable level, the Verlinde formula
\eqref{eq:verlinde-general} provides an independent
computation; the problem is to derive the Verlinde formula
directly from the Feynman integral expansion
\eqref{eq:thqg-X-class-L-correction}.
\end{openproblem}

\begin{openproblem}[Analytic completion for W-algebras]%
% label removed: prob:thqg-X-W-algebra-analytic%
\index{W-algebra!analytic completion}%
Construct the sewing envelope and analytic bar coalgebra for
class-M algebras ($\mathrm{Vir}_c$, $\cW_N$) in the
resonance-filtered topology of
Theorem~\ref{thm:resonance-filtered-bar-cobar}.
The resonance rank $\rho(\cA)$
(Definition~\ref{def:resonance-rank}) controls the
difficulty: for $\rho < \infty$ (the resonance completion
theorem, Theorem~\ref{thm:platonic-completion}), the
analytic completion reduces to a finite resonance problem;
for $\rho = \infty$ (if such algebras exist among modular
Koszul chiral algebras), the completion requires a genuinely
infinite-dimensional analytic framework beyond the Fredholm
theory of this section.
\end{openproblem}

\begin{remark}[Summary: the Fredholm partition function as
a projection of $\Theta_\cA$]%
% label removed: rem:thqg-X-fredholm-summary%
\index{Fredholm determinant!as MC projection}%
The genus-$g$ Fredholm partition function
$Z_g(\cA;\,\Omega) = \det(1 - \sewker_g)^{-\kappa(\cA)}$
is the analytic realization of the scalar genus tower
$F_g(\cA) = \kappa(\cA) \cdot \lambda_g^{\mathrm{FP}}$
of Theorem~D.  It is the bottom projection of the single
universal MC element
$\Theta_\cA \in \MC(\gAmod)$
through the analytic completion:
\[
\Theta_\cA
\;\xrightarrow{\;\mathrm{arity}\;2\;}
\kappa(\cA)
\;\xrightarrow{\;\mathrm{sewing}\;}
\sewker_g
\;\xrightarrow{\;\det\;}
Z_g.
\]
For class-G algebras, this projection is lossless: $Z_g$
captures all the information in~$\Theta_\cA$.  For classes
L, C, M, the non-Fredholm corrections $\delta Z_g$ carry the
information from the higher-arity projections of~$\Theta_\cA$
(the cubic shadow $\fC$, the quartic resonance class
$\mathfrak{Q}$, and all higher obstruction classes $o_r$).
The full partition function
$Z_g^{\mathrm{full}} = Z_g^{\mathrm{G}} + \delta Z_g$
is the analytic realization of the complete MC element, and
the Feynman expansion \eqref{eq:thqg-X-feynman-expansion}
is the graph-sum formula of the shadow obstruction tower
evaluated on the analytic sewing data.

This completes the development of Result~\textbf{(G10)}: the
genus-$g$ partition function of a modular Koszul chiral algebra,
computed as a nonperturbative Fredholm determinant on the
one-particle Bergman space, with explicit Feynman corrections
for non-Gaussian shadow obstruction towers.
\end{remark}

%% ----------------------------------------------------------------
%% Analytic Borel--Moore homology and Fredholm Koszul duality
%% ----------------------------------------------------------------

\subsection{Analytic Borel--Moore homology and Fredholm Koszul duality}
% label removed: subsec:analytic-bm-steinberg

\providecommand{\fSb}{\mathfrak{S}_b}

The Steinberg variety $\fSb$ attached to a modular Koszul chiral
algebra $\cA$ parametrises composable pairs of bar elements:
a point of~$\fSb$ is a pair $(b_1, b_2)$ in
$\anbar(\cA) \times_{\cA} \anbar(\cA)$ meeting along a common
$\cA$-output.  The convolution product on~$\fSb$ encodes
the bar differential in the holomorphic direction and the
deconcatenation coproduct in the topological direction,
so that the resulting algebra is the Fredholm-level avatar
of the Swiss-cheese operad action.

\begin{construction}[Analytic Borel--Moore homology of $\fSb$]%
% label removed: constr:analytic-bm-steinberg%
\index{Steinberg variety!analytic BM homology}%
\index{Borel--Moore homology!analytic}%
Complete the Steinberg object $\fSb$ in the analytic topology
appropriate to the shadow class of~$\cA$:
\begin{itemize}
\item \emph{Class-F} (free field): the nuclear Fr\'echet topology
  on the Bergman one-particle space, so that the FM strata
  of~$\fSb$ carry the topology of nuclear Fr\'echet--Stein
  algebras.
\item \emph{Class-M} (W-algebras, resonance): the
  resonance-filtered topology of
  Theorem~\ref{thm:resonance-filtered-bar-cobar}, in which
  $\fSb$ is exhausted by Banach-analytic subvarieties of
  finite resonance rank~$\rho(\cA)$.
\end{itemize}
In either case, define the \emph{analytic Borel--Moore homology}
$H_*^{\mathrm{BM,an}}(\fSb)$ as the homology of the complex of
locally finite analytic chains on~$\fSb$ with coefficients in
the dualising sheaf.  The convolution product
\[
H_i^{\mathrm{BM,an}}(\fSb)
\otimes
H_j^{\mathrm{BM,an}}(\fSb)
\;\xrightarrow{\;\mu\;}
H_{i+j}^{\mathrm{BM,an}}(\fSb)
\]
makes $H_*^{\mathrm{BM,an}}(\fSb)$ into an associative graded
algebra, the \emph{Fredholm-level convolution algebra} of~$\cA$.
\end{construction}

\begin{conjecture}[\ClaimStatusConjectured]%
% label removed: prop:analytic-steinberg-koszul%
\index{Koszul duality!Fredholm-level Steinberg}%
\index{coderived category}%
\index{contraderived category}%
For a modular Koszul chiral algebra $\cA$ satisfying HS-sewing,
there is an equivalence of triangulated categories
\begin{equation}% label removed: eq:analytic-steinberg-koszul
\codercat\bigl(H_*^{\mathrm{BM,an}}(\fSb)\textup{-comod}\bigr)
\;\simeq\;
\contrcat(\cA^!\textup{-mod}).
\end{equation}
This is the analytic enhancement of the algebraic Steinberg
presentation (Theorem~\ref{thm:steinberg-presentation}):
$H_*^{\mathrm{BM,an}}(\fSb)$ presents the Fredholm-level
bar coalgebra, and
\eqref{eq:analytic-steinberg-koszul} is the analytic lift
of bar-cobar inversion (Theorem~B).
The coderived/contraderived asymmetry reflects the choice
of completion topology on~$\fSb$: the coderived side
carries the direct-limit (nuclear) topology of bar comodules,
while the contraderived side carries the inverse-limit
(pro-finite-dimensional) topology of Koszul dual modules.
\end{conjecture}

\begin{remark}[Non-Fredholm corrections from strata at infinity]%
% label removed: rem:steinberg-non-fredholm%
\index{non-Fredholm correction!strata at infinity}%
\index{Siegel modular form}%
For class-L algebras, the non-Fredholm correction
$\delta Z_g$ of
\S\ref{subsec:thqg-non-fredholm} receives a Steinberg
interpretation: it is the contribution from the
\emph{non-compact strata} of~$\fSb$, the strata at
infinity in the FM compactification
$\overline{\mathrm{FM}}_k(\Sigma_g)$.  The genus-$g$
Steinberg variety carries a natural fibration
\[
\fSb^{(g)} \;\longrightarrow\; \mathbb{H}_g
\]
over the Siegel upper half-space $\mathbb{H}_g$, where
the fibre over $\Omega \in \mathbb{H}_g$ is the Steinberg
variety of the sewing kernel
$\sewker_g(\Omega)$.  The compact strata produce the
Fredholm determinant $Z_g^{\mathrm{G}}$; the non-compact
strata produce
\[
\delta Z_g^{\mathrm{L}}
\;=\;
\int_{\mathbb{H}_g \smallsetminus \mathcal{F}_g}
\operatorname{ch}\bigl(
  H_*^{\mathrm{BM,an}}(\fSb^{(g)}_\Omega)
\bigr)
\,d\mu_{\mathrm{Siegel}},
\]
where $\mathcal{F}_g \subset \mathbb{H}_g$ is a Siegel
fundamental domain and $\operatorname{ch}$ denotes the
Chern character of the convolution algebra fibre.  At
integrable level, this integral localises to the
Verlinde formula via the finite-dimensionality of the
conformal-block bundle.
\end{remark}

% ============================================================
\subsection{Prime-locality and quantum gravity}%
\label{subsec:prime-locality-quantum-gravity}%
\index{prime-locality!quantum gravity interface|textbf}

The prime-locality programme
(Vol~I, Conjecture~conj:prime-locality-transfer and
\S\ref{sec:prime-locality-assessment})
asks whether the shadow amplitudes at genus~$1$ decompose
into Hecke eigencomponents with multiplicative Fourier
coefficients.  The question has a precise translation into the
$3$d quantum gravity framework of this volume: prime-locality
is the question of whether the holomorphic-topological partition
function, viewed on the genus-$1$ moduli space, respects the
arithmetic structure of the Hecke correspondences.

\begin{remark}[What the $3$d perspective adds to prime-locality]%
\label{rem:3d-prime-locality}%
\index{prime-locality!3d perspective}
The $3$d holomorphic-topological framework provides three
structural inputs that bear on the prime-locality question:

\textbf{Input~1: the sewing-Hecke collapse is a genus-$1$
phenomenon.}
The sewing operator $\sewop_q$ at genus~$1$ is diagonal in the
Fock basis: $\sewop_q|n\rangle = q^n|n\rangle$.  This
diagonality forces the sewing--Hecke reciprocity
(Vol~I, Theorem~thm:sewing-hecke-reciprocity), which collapses
the two-variable spectral data into a single-variable Dirichlet
series.  At genus~$\ge 2$, the off-diagonal entries $R_{ij}$
in the multi-handle sewing kernel
(\S\ref{subsec:thqg-genus2-fredholm}) break this diagonality.
Consequently, the arithmetic content at genus~$\ge 2$ is
genuinely multi-variable: the Rankin--Selberg unfolding on
$\mathrm{Sp}_{2g}(\mathbb{Z}) \backslash \mathfrak{H}_g$
produces Siegel modular $L$-functions that do not factor
as products of elliptic $L$-functions.

\textbf{Consequence:} the prime-locality question is tied to
the genus-$1$ sewing operator.  At higher genus, the relevant
arithmetic structure is not prime-locality but
\emph{Siegel--Hecke equivariance}, a genuinely different (and
harder) problem.  The genus-$2$ B\"ocherer bridge
(Vol~I, genus-$2$ arithmetic frontier) is the first probe
of this regime.

\textbf{Input~2: the gravitational dichotomy constrains the
Hecke defect.}
At $c = 26$, the matter-ghost system has $\kappa_{\mathrm{eff}}
= \kappa(\mathrm{matter}) + \kappa(\mathrm{ghost}) = 0$,
and the transgression algebra degenerates from Clifford to
exterior (Vol~I, the critical string dichotomy).  At
$c \neq 26$, the transgression algebra is a matrix algebra
(Morita-trivial), and the Koszul duality pairing
$\cA \leftrightarrow \cA^!$ is non-degenerate.

For the Hecke defect class
$[\delta_p] \in H^1(\mathfrak{g}_\cA^{\mathrm{mod}},
\mathrm{ad}_{T_p})$,
the Morita triviality at $c \neq 26$ means that the
convolution algebra has ``fewer'' cohomological directions
for the defect to live in.  At $c = 26$, the degeneration
of the transgression algebra opens new cohomological
directions.  Whether this structural difference affects the
vanishing of $[\delta_p]$ is a concrete question that
connects prime-locality to the critical string dichotomy.

\textbf{Input~3: the planted-forest correction is
tau-dependent at genus~$\ge 2$.}
The planted-forest differential $d_{\mathrm{pf}}$ of
Vol~I involves codimension-$\ge 2$ strata of the log-FM
compactification.  At genus~$1$, the planted-forest
correction $\delta_{\mathrm{pf}}^{(1,r)}$ contributes to the
shadow amplitudes but does not introduce new tau-dependence
beyond $E_2^*(\tau)$.  At genus~$\ge 2$, the planted-forest
corrections involve Siegel modular functions, and the
Hecke equivariance question becomes the question of whether
the planted-forest graph sums on $\overline{\mathcal{M}}_{g,n}$
respect the Siegel--Hecke decomposition.
\end{remark}

\begin{remark}[The honest assessment from the $3$d perspective]%
\label{rem:3d-prime-locality-assessment}%
\index{prime-locality!honest assessment}
The $3$d holomorphic-topological framework does \emph{not}
resolve the prime-locality question.  It clarifies the
\emph{landscape}:
\begin{enumerate}[label=\textup{(\roman*)}]
\item Prime-locality is proved for lattice VOAs
  (all ranks, all levels) and for rational VOAs with diagonal
  modular invariant.  The proof uses the Hecke module structure
  of the theta function or the Franc--Mason eigendecomposition,
  not the $3$d framework.

\item Prime-locality at the shadow \emph{coefficient} level
  (tau-independent OPE data) is proved for all algebraic families
  by tau-independence alone.  This is a trivial consequence of
  the fact that the Hecke operator acts as the identity on
  constants.

\item The open question is prime-locality at the shadow
  \emph{amplitude} level for irrational VOAs: whether the
  genus-$1$ graph sums
  $\operatorname{Sh}_r^{(1)}(\tau) =
  \sum_\Gamma |\mathrm{Aut}(\Gamma)|^{-1}
  \prod_{v} S_{n_v} \prod_{e} E_2^*(\tau)$
  respect the Hecke decomposition of the ambient space of
  quasi-modular forms.  The vertex weights $S_{n_v}$ are
  Hecke-equivariant (tau-independent); the edge weights
  $E_2^*(\tau)$ are individually Hecke-equivariant (the
  Fourier coefficients $\sigma_1(n)$ are multiplicative); but
  the \emph{products} of $E_2^*$ are not automatically
  Hecke-equivariant, because $T_p$ is not a ring homomorphism
  on quasi-modular forms.

\item The $3$d framework contributes one genuine insight:
  the prime-locality question is \emph{intrinsically genus-$1$}.
  At higher genus, the relevant arithmetic object is not a
  Hecke eigendecomposition on $\mathcal{M}_{1,1}$ but a
  Siegel--Hecke structure on $\mathcal{M}_g$, and the
  sewing-kernel non-diagonality opens different
  (and potentially richer) arithmetic channels.
\end{enumerate}

The Riccati algebraicity theorem
(Vol~I, Theorem~thm:riccati-determinacy-assessment)
provides a well-defined ``algebraic prime-locality'': the
shadow obstruction tower is finitely determined by three OPE invariants
$(\kappa, \alpha, S_4)$.  This is a strong constraint on the
algebraic structure but does not address the arithmetic question
of Hecke eigenvalues.  The term ``prime-locality'' should be
reserved for the arithmetic sense; the algebraic determinacy
is better called \emph{Riccati determinacy} or
\emph{shadow-tower finite determination}.
\end{remark}

\begin{remark}[Siegel modular forms and the genus-$2$ HT partition function]%
\label{rem:thqg-siegel-genus2}%
\index{Siegel modular form!3d HT partition function}%
\index{Fredholm determinant!genus-2}%
The genus-$2$ Fredholm determinant
$\freddet(1 - K_2)$ of
Proposition~\textup{thm:genus2-non-diagonal}
\textup{(}Vol~I\textup{)} produces a function on
$\Siegel_2 = \Sp(4,\bZ) \backslash \mathfrak{H}_2$
whose Fourier expansion connects to the Igusa ring
$M_*(\Sp(4,\bZ)) = \bC[E_4, E_6, \chi_{10}, \chi_{12}]$
of Siegel modular forms.
For lattice VOAs of rank~$d$, the genus-$2$ partition
function is the Siegel theta series
$\Theta_\Lambda^{(2)}(\Omega)$, a Siegel modular form of
weight~$d/2$.  The three-dimensional HT perspective adds
a structural interpretation: the off-diagonal coupling
$R_{12}$ in the Schur complement factorization is the
genus-$2$ manifestation of the bar complex coproduct
\textup{(}the $\bR$-direction factorization of the
Swiss-cheese structure\textup{)}, and the cusp-form
component of $\Theta_\Lambda^{(2)}$ encodes
genuinely genus-$2$ bulk--boundary correlations that
are invisible to the genus-$1$ shadow obstruction tower.

For the Leech lattice \textup{(}$d = 24$\textup{)}:
the cusp projection onto the Saito--Kurokawa lift
$\chi_{12} = \mathrm{SK}(f_{22})$ is nonzero
\textup{(}Vol~I, Proposition~\textup{prop:leech-cusp-nonvanishing}\textup{)},
and the B\"ocherer coefficient provides algebraic access
to the central $L$-value $L(1/2, \pi_{f_{22}} \times \chi_D)$
\textup{(}Vol~I, Theorem~\textup{thm:bocherer-bridge}\textup{)}.
In the 3d HT framework, this is the genus-$2$ version of
the holographic modular Koszul datum:
the genus-$2$ MC equation determines the Siegel theta
series from genus-$1$ shadow data and boundary
contributions, and the B\"ocherer factorization converts
this algebraic determination into critical-line access.
\end{remark}

\begin{remark}[MZV periods and the genus-$0$ bar complex
in 3d HT]%
\label{rem:thqg-mzv-bar}%
\index{multiple zeta values!3d HT}%
\index{bar complex!MZV periods in 3d HT}%
The genus-$0$ bar amplitudes of the holomorphic colour
produce multiple zeta value (MZV) periods via iterated
integrals on $\cM_{0,n}$
\textup{(}Vol~I, Theorem~\textup{thm:shadow-mzv-dictionary}\textup{)}.
In the 3d HT interpretation: the MZV periods are the
perturbative contributions to the holomorphic Wilson line,
computed by the bar complex on $\overline{\mathrm{FM}}_n(\bC)$.
The shadow depth classification $\mathbf{G}/\mathbf{L}/\mathbf{C}/\mathbf{M}$
controls which MZV weights appear:
\begin{enumerate}[label=\textup{(\roman*)}]
\item For 3d HT theories with Heisenberg boundary
\textup{(}class~$\mathbf{G}$\textup{)}: only $\zeta(2)$
appears, and the perturbative holomorphic Wilson line
terminates at one loop.
\item For 3d HT theories with affine Kac--Moody boundary
\textup{(}class~$\mathbf{L}$\textup{)}: $\zeta(2)$ and
$\zeta(3)$ appear.  The Drinfeld associator
\textup{(}Vol~I,
Theorem~\textup{thm:drinfeld-associator-bar-transport}\textup{)}
is the parallel transport of the genus-$0$ bar
connection, and its MZV coefficients encode the
perturbative holomorphic data.
\item For 3d gravity \textup{(}Virasoro boundary,
class~$\mathbf{M}$\textup{)}: MZVs at all weights
appear, reflecting the infinite tower of perturbative
corrections to the holomorphic graviton exchange.
The Zagier--Brown dimension sequence
$d_r = d_{r-2} + d_{r-3}$ counts the independent
perturbative MZV contributions at each order.
\end{enumerate}
\end{remark}
